\documentclass[DIV=20,oneside]{scrartcl}
\usepackage{graphicx}
\usepackage{amsmath}
\usepackage{float}
\usepackage{enumitem}

\usepackage{scrlayer-scrpage}

\clearpairofpagestyles
\ihead{Dr. Srikantha N.}
\setheadsepline{0.4pt}
\pagestyle{scrheadings}

\usepackage{url}
\usepackage{cancel}

\usepackage[most]{tcolorbox}

\usepackage{listings}
\usepackage{xcolor}
\lstset{
	basicstyle=\ttfamily\small,
	frame=single,
	backgroundcolor=\color{gray!10},
	keywordstyle=\color{blue},
	commentstyle=\color{green!60!black},
	numbers=left,
	numberstyle=\tiny\color{gray},
	breaklines=true,
	breakatwhitespace=true,
	columns=fullflexible,
	keepspaces=true
}

\newcounter{problem}
\newcommand{\problem}{
	\stepcounter{problem}
	\noindent\textbf{\\Problem \theproblem: }
}

\newcommand{\note}{
	\noindent\textbf{\\Note: }
}

\newcommand{\insight}[1]{%
	\begin{tcolorbox}[
		colback=black!20,
		colframe=black,
		boxrule=0.5pt,
		arc=2mm]
		\textbf{Insight: } #1
	\end{tcolorbox}
}

\newcommand{\takeaway}[1]{%
	\begin{tcolorbox}[
		colback=black!20,
		colframe=black,
		boxrule=0.5pt,
		arc=2mm]
		\textbf{Takeaway: } #1
	\end{tcolorbox}
}

\newcommand{\solution}{
	\noindent\textbf{\\Solution: }
}
\newcommand{\example}{
	\noindent\textbf{\\Example: }
}
\newcommand{\proof}{
	\noindent\textbf{\\Proof: }
}
\newcommand{\deff}{
	\noindent\textbf{\\Definition: }
}
%opening
\title{Linear Algebra and Optimization }
\author{Dr. Srikantha N, RIT, Bengaluru}

\begin{document}
	\date{}
	\maketitle
	
	\section{Matrices, Linear Equations and Vector Space:}
	\deff
	A vector is an ordered sequence or list of numbers. \textbf{Example: } $v=[1,2,3]$ or $v = \begin{bmatrix} 1 \\ 2 \\ 3 \end{bmatrix}$
	\par In the context of Data Science and Machine Learning, we interpret a vector as a feature vector representing a single data instance or sample. If we collect $d$ distinct attributes (features) about an object, that object is represented as a vector residing in a $d$-dimensional space, denoted as ${R}^d$. Mathematical Representation: A vector $\mathbf{x} \in {R}^d$ is typically expressed as a column matrix:$$\mathbf{x} = \begin{bmatrix} x_1 \\ x_2 \\ \vdots \\ x_d \end{bmatrix}$$
	\example: \textbf{Text and Natural Language Processing (NLP):}
	In document analysis or search engines, a piece of text can be mathematically encoded into a vector using term frequencies (or embeddings). Let's represent a short document based on the vocabulary count of five words:
	\\
	\textbf{Text: }"Data security is important because data leaks damage trust."
	\\
	\textbf{Feature Space: }\texttt{["data", "security", "leaks", "trust", "algorithm"]}
	\\
	\textbf{Resulting Vector $\mathbf{x_{\text{doc}}}$:}
	$$\mathbf{x}_{\text{doc}} = \begin{bmatrix} 2 \\ 1 \\ 1 \\ 1 \\ 0 \end{bmatrix}$$
	\insight{In NLP, Vocabulary (The Vector Space) is a master list of all unique words the AI is programmed to recognize. Sentiment Weight Vector is a matching vector of numbers that holds the emotional score for every single word in the vocabulary. These weights aren't guessed; they are learned by the AI during training after reading millions of pre-labeled reviews (like 5-star or 1-star ratings). Words that appear constantly in positive reviews get high positive weights, while words frequent in complaints get negative weights. Neutral grammar words sit at zero.}
	\example Now consider a vocabulary set \texttt{['and', 'also', 'but', 'difficult', 'easy', 'good', 'great', 'is', 'not']}\\
	Weights for the words in vocabulary:
	\begin{itemize}[itemsep=0.5pt, topsep=1pt]
		\item Positive: great (+3.0), easy (+2.5), good (+2.0)
		\item Negative: difficult (-2.0), but (-0.5), not (-0.5)
		\item Neutral: and (0.0), also (0.0), is (0.0)
	\end{itemize}
	
	$$\mathbf{w} = [0.0, 0.0, -0.5, -2.0, 2.5, 2.0, 3.0, 0.0, -0.5]$$
	Now consider 2 sentences:
	\begin{center}
		Sentence 1: "NLP is great and also easy"\\
		Sentence 2: "NLP is not good but also difficult"
	\end{center}
	Sentence 1: Vector ($\mathbf{x}_1$): [1, 1, 0, 0, 1, 0, 1, 1, 0]\quad \quad 	Dot Product Score: $(1 \times 2.5) + (1 \times 3.0) = \mathbf{5.5}$, AI Label: Overwhelmingly Positive review\\
	Sentence 2: Vector ($\mathbf{x}_2$): [0, 1, 1, 1, 0, 1, 0, 1, 1]\quad \quad 
	Dot Product Score: $(-0.5) + (-2.0) + (1 \times 2.0) + (-0.5) = \mathbf{-1.0}$,
	AI Label: Strongly Negative review
	
	\takeaway{\textbf{Word2vec} is a technique in natural language processing for obtaining vector representations of words. Typically, the dimensionality of the vectors is set to be between 100 and 1,000.\textbf{OpenAI’s text embeddings} measure the relatedness of text strings. By default, the length of the embedding vector is 1536 for text-embedding-3-small or 3072 for text-embedding-3-large. The hidden internal hidden-state embedding size of the core 175B \textbf{GPT-3 model} architecture is 12,288 dimensions. }
	
	\example \textbf{Image Analytics:} An image can be mathematically represented as a 2D matrix of pixel intensity values. When we prepare this image for a machine learning model or computational analysis, we flatten this 2D matrix into a 1D vector.Here is the representation of a $4 \times 4$ grayscale image filled with intensity values of $0$ (black), $0.5$ (gray), and $1$ (white).
	\\
	Arranged as a grid (where rows and columns map to the spatial layout of the pixels), the matrix $M$ looks like this:$$M = \begin{bmatrix} 
		0.0 & 0.5 & 0.5 & 0.0 \\ 
		0.5 & 1.0 & 1.0 & 0.5 \\ 
		0.5 & 1.0 & 1.0 & 0.5 \\ 
		0.0 & 0.5 & 0.5 & 0.0 
	\end{bmatrix}$$
	\begin{figure}[H]
		\centering
		\includegraphics[width=0.3\linewidth]{BinaryImage}
		\caption{Digital Image}
		\label{fig:binaryimage}
	\end{figure}
		To convert this image into data suitable for a neural network or vector operations, we unroll or "flatten" the matrix row-by-row (known as row-major order). This transforms the $4 \times 4$ spatial grid into a 16-dimensional feature vector $\mathbf{x} \in {R}^{16}$:
		\[
		\mathbf{x}^T=
		\left[
		\begin{array}{cccccccccccccccc}
			0.0 & 0.5 & 0.5 & 0.0 & 0.5 & 1.0 & 1.0 & 0.5 &
			0.5 & 1.0 & 1.0 & 0.5 & 0.0 & 0.5 & 0.5 & 0.0
		\end{array}
		\right]
		\]
		Now by replacing $0.0 \to 1.0$, $1.0 \to 0.0$, and keeping $0.5 \to 0.5$, we get the following spatial matrix representation:
		$$M_{\text{inverted}} = \begin{bmatrix} 
			1.0 & 0.5 & 0.5 & 1.0 \\ 
			0.5 & 0.0 & 0.0 & 0.5 \\ 
			0.5 & 0.0 & 0.0 & 0.5 \\ 
			1.0 & 0.5 & 0.5 & 1.0 
		\end{bmatrix}$$
		and 
		\[
		\mathbf{y}^T=
		\left[
		\begin{array}{cccccccccccccccc}
			1.0 & 0.5 & 0.5 & 1.0 & 0.5 & 0.0 & 0.0 & 0.5 &
			0.5 & 0.0 & 0.0 & 0.5 & 1.0 & 0.5 & 0.5 & 1.0
		\end{array}
		\right]
		\]
		Corresponding image is as follows:
		\begin{figure}[H]
			\centering
			\includegraphics[width=0.5\linewidth]{CosineSimilarity1}
			\caption{Similarity Search Example 01}
			\label{fig:cosinesimilarity1}
		\end{figure}
		\subsection{Cosine Similarity: }Cosine Similarity measures the cosine of the angle between two multi-dimensional vectors. It evaluates directional orientation rather than scale, ranging from $-1$ (exactly opposite) to $+1$ (exactly identical).
		$$\text{Cosine Similarity} = \frac{\mathbf{x} \cdot \mathbf{y}}{\|\mathbf{x}\| \|\mathbf{y}\|}$$
		For the above example we get
		$$\text{Cosine Similarity} = \frac{2.0}{\sqrt{6} \times \sqrt{6}} = \frac{2.0}{6.0} = \frac{1}{3} \approx 0.3333$$
		A score of $\frac{1}{3}$ shows a mild positive similarity. Even though we flipped the structural extreme opposites (black and white pixels), they still share a common, un-flipped spatial feature framework—namely, the gray borders ($0.5$), which preserve enough shared signal weight to keep the similarity metric above $0.0$.
		\\ If we had completely inverted the gray values as well (e.g., $0.5 \to -0.5$ or mapping to a strict binary inverse), the dot product would collapse further, signaling a negative or orthogonal relational correlation.
		
		Let's modify the image vector by changing only the right corners from black ($0.0$) to white ($1.0$), while keeping all other pixels identical to the original image $\mathbf{x}$. We get:
		
		\begin{figure}[H]
			\centering
			\includegraphics[width=0.5\linewidth]{CosineSimilarity}
			\caption{Similarity Search Example 02}
			\label{fig:cosinesimilarity}
		\end{figure}
		$$\text{Cosine Similarity} = \frac{6.0}{\sqrt{6} \times \sqrt{8}} = \frac{6.0}{\sqrt{48}} = \frac{6.0}{4\sqrt{3}} = \frac{3}{2\sqrt{3}} = \frac{\sqrt{3}}{2} \approx 0.8660$$
		Flipping just two out of the sixteen pixels from black to white yields a high similarity score of $86.6\%$.
		\\
		\subsection{System of Linear Equations:} A collection of one or more linear equations involving the same set of variables. It can be written compactly in matrix form as:$$Ax = b$$
		Where:
		\begin{itemize}[itemsep=0.5pt, topsep=1pt]
			\item $A$ is an $m \times n$ matrix of coefficients (e.g., feature matrix).
			\item $x$ is an $n \times 1$ vector of unknowns (e.g., model weights or parameters).
			\item $b$ is an $m \times 1$ output vector (e.g., target labels or ground truth).
		\end{itemize}
		\example \textbf{1:}
		$$\begin{aligned}
			1x_1 + 2x_2 &= 5 \\
			3x_1 + 4x_2 &= 11
		\end{aligned}$$
		This system can be written in the matrix form $Ax = b$ as:
		$$\begin{bmatrix} 1 & 2 \\ 3 & 4 \end{bmatrix} \begin{bmatrix} x_1 \\ x_2 \end{bmatrix} = \begin{bmatrix} 5 \\ 11 \end{bmatrix}$$
		\example \textbf{2}:
		Let's look at a concrete dataset of house pricing: \\Feature 1 ($x_1$): Size of the house (in thousands of square feet). \\
		Feature 2 ($x_2$): Number of bedrooms.Target ($b$): 
		Price of the house (in hundreds of thousands of dollars). 
		\\
		Suppose we have gathered data from two houses:\\
		House 1: 1,000 sq. ft ($x_1 = 1$), 2 bedrooms ($x_2 = 2$), sold for Rs. 5 lakhs ($b_1 = 5$).\\
		House 2: 3,000 sq. ft ($x_1 = 3$), 4 bedrooms ($x_2 = 4$), sold for Rs. 11 lakhs ($b_2 = 11$).
		\\
		We want to find the model weights ($w_1, w_2$) that satisfy this data exactly:$$\begin{bmatrix} 1 & 2 \\ 3 & 4 \end{bmatrix} \begin{bmatrix} w_1 \\ w_2 \end{bmatrix} = \begin{bmatrix} 5 \\ 11 \end{bmatrix}$$By solving this system (using methods like Gauss-Jordan elimination), we find $w_1 = 1$ and $w_2 = 2$. This gives us our predictive model:
		$$\text{Price} = 1 \times (\text{Size}) + 2 \times (\text{Bedrooms})$$
		\subsubsection{Geometrical Interpretation: }
			Geometrically, each linear equation in a system represents a hyperplane in an $n$-dimensional space ${R}^n$. In a 2D space (${R}^2$), a linear equation represents a line. In a 3D space (${R}^3$), a linear equation represents a flat plane. The solution to a system of linear equations is the intersection of these hyperplanes. There are three possible geometric outcomes for a linear system: 
			\\
			\textbf{Unique Solution:} The hyperplanes intersect at exactly one point. (The system is consistent and independent).
			$$\text{rank}(A) = \text{rank}([A | b]) = n$$
			\\
			\textbf{Infinite Solutions:} The hyperplanes align perfectly or intersect along an entire line/subspace, meaning they represent redundant features. (The system is consistent and dependent).
			$$\text{rank}(A) = \text{rank}([A | b]) < n$$
			\\
			\textbf{No Solution:} The hyperplanes are parallel and never intersect, meaning the target vector $b$ lies outside the column space of $A$. (The system is inconsistent).
			$$\text{rank}(A) < \text{rank}([A | b])$$
			\\
			Let us visually contrast a unique solution system against a system with no solution using lines in ${R}^2$: \\
			Case A (Unique Solution):$$\begin{aligned}
				x_1 + x_2 &= 3 \\
				x_1 - x_2 &= 1
			\end{aligned}$$These represent two distinct lines that intersect exactly at the point $(2, 1)$.\\
			Case B (No Solution):$$\begin{aligned}
				x_1 + x_2 &= 3 \\
				2x_1 + 2x_2 &= 8
			\end{aligned}$$ These represent two parallel lines that will never cross.
			\begin{figure}[H]
				\centering
				\includegraphics[width=0.7\linewidth]{GeometricSolution}
				\caption{Geometrical interpretation of the solution to system of equations}
				\label{fig:geometricsolution}
			\end{figure}
	\problem
	Solve the following $3 \times 3$ linear system and interpret it geometrically:$$\begin{aligned}
		x_1 + x_2 + x_3 &= 6 \\
		2x_1 + 3x_2 + 2x_3 &= 14 \\
		x_1 + 2x_2 + 3x_3 &= 14
	\end{aligned}$$
	$$[A | b]=\begin{bmatrix} 1 & 1 & 1 & | & 6 \\ 2 & 3 & 2 & | & 14 \\ 1 & 2 & 3 & | & 14 \end{bmatrix}$$		
	Using these row operations:
	$R_2 \leftarrow R_2 - 2R_1$:$R_3 \leftarrow R_3 - R_1$:$R_3 \leftarrow R_3 - R_2$: we get:
	$$[A|b]=\begin{bmatrix} 1 & 1 & 1 & | & 6 \\ 0 & 1 & 0 & | & 2 \\ 0 & 0 & 2 & | & 6 \end{bmatrix}$$
	$$\text{rank}(A) = \text{rank}([A | b]) = n$$
	The unique solution vector is $x = \begin{bmatrix} 1 , 2 , 3 \end{bmatrix}$
	\problem Solve the following $3 \times 3$ linear system and interpret it geometrically:$$\begin{aligned}
		x_1 + 2x_2 + x_3 &= 4 \\
		2x_1 + 5x_2 + 4x_3 &= 11 \\
		3x_1 + 7x_2 + 5x_3 &= 15
	\end{aligned}$$
	$$[A | b]=\begin{bmatrix} 1 & 2 & 1 & | & 4 \\ 2 & 5 & 4 & | & 11 \\ 3 & 7 & 5 & | & 15 \end{bmatrix}$$
	Performing these row operations: $R_2 \leftarrow R_2 - 2R_1$:$R_3 \leftarrow R_3 - 3R_1$:$R_3 \leftarrow R_3 - R_2$: we get:
	$$[A|b]=\begin{bmatrix} 1 & 2 & 1 & | & 4 \\ 0 & 1 & 2 & | & 3 \\ 0 & 0 & 0 & | & 0 \end{bmatrix}$$
	$$\text{rank}(A) = \text{rank}([A | b]) < n$$
	Hence infinitely many solutions. 
	This gives these following system of equations: $x_1+2x_2+x_3=4: x_2+2x_3=3$\\
	Let $x_3$ be a free parameter, $x_3 = t\implies x_2 + 2t = 3 \Rightarrow x_2 = 3 - 2t\implies  x_1 = -2 + 3t$.
	
	$$x = \begin{bmatrix} -2 + 3t , 3 - 2t , t \end{bmatrix}$$
	\problem
	Solve the following $3 \times 3$ linear system and interpret it geometrically:$$\begin{aligned}
		x_1 + x_2 + 2x_3 &= 3 \\
		2x_1 + 3x_2 + 5x_3 &= 7 \\
		3x_1 + 4x_2 + 7x_3 &= 12
	\end{aligned}$$
	$$[A | b]=\begin{bmatrix} 1 & 1 & 2 & | & 3 \\ 2 & 3 & 5 & | & 7 \\ 3 & 4 & 7 & | & 12 \end{bmatrix}=\begin{bmatrix} 1 & 1 & 2 & | & 3 \\ 0 & 1 & 1 & | & 1 \\ 0 & 0 & 0 & | & 2 \end{bmatrix}$$
	$$\text{rank}(A) < \text{rank}([A | b])$$
	Consequently, the system yields no solution.
	\subsection{Gauss-Jordan method}
	While standard Gaussian elimination stops at an Upper Triangular Matrix (Row Echelon Form - REF) and requires back-substitution to find the variables, the Gauss-Jordan method goes a step further. It uses elementary row operations to reduce a matrix completely into Reduced Row Echelon Form (RREF), where the solution can be read directly from the matrix.
	\\A matrix is in Reduced Row Echelon Form (RREF) if it satisfies these conditions:
	\begin{itemize}[itemsep=0.5pt, topsep=1pt]
		\item All rows consisting entirely of zeros are at the bottom of the matrix.
		\item The first non-zero entry in any row is a 1 (called a leading 1 or a pivot).
		\item Each leading 1 is the only non-zero entry in its entire column (meaning zeros must be both below and above it).
		\item 	The leading 1 of a lower row occurs strictly to the right of the leading 1 of a higher row.
	\end{itemize}
	\example Solve the following system using the Gauss-Jordan method:$$\begin{aligned}
		x + y + z &= 6 \\
		2y + 5z &= -4 \\
		2x + 5y - z &= 27
	\end{aligned}$$
	$$[A | b]=\begin{bmatrix} 
		1 & 1 & 1 & | & 6 \\ 
		0 & 2 & 5 & | & -4 \\ 
		2 & 5 & -1 & | & 27 
	\end{bmatrix}$$
	Using these row operations: $R_3 \leftarrow R_3 - 2R_1$: $R_3 \leftarrow 2R_3-3R_2$: $R_2 \leftarrow 21R_2+5R_3$: $R_1 \leftarrow 21R_1+R_3$: $R_1 \leftarrow 2R_1-R_2:$
	$$[A|b]=\begin{bmatrix} 
		42 & 0 & 0 & | & 210 \\ 
		0 & 42 & 0 & | & 126 \\ 
		0 & 0 & -21 & | & 42 
	\end{bmatrix}=\begin{bmatrix} 
	1 & 0 & 0 & | & 5 \\ 
	0 & 1 & 0 & | & 3 \\ 
	0 & 0 & 1 & | & -2 
	\end{bmatrix}$$
	Final Solution: $x = 5, \quad y = 3, \quad z = -2$
	\subsection{LU Decomposition: }
	The LU decomposition is a process of factorizing an $n \times n$ data or coefficient matrix $A$ into two triangular matrices:
	$$A = LU$$
	Where: 
	\begin{itemize}[itemsep=0.5pt, topsep=1pt]
		\item $L$ is a Lower Triangular Matrix (all entries above the main diagonal are $0$, and the diagonal entries are typically $1$).
		\item $U$ is an Upper Triangular Matrix (all entries below the main diagonal are $0$).
	\end{itemize}
	\insight{{\textbf{\textit{If we can just compute $A^{-1}$ using Gauss Elimination and Gauss Jordan, why do we need LU factorization?}}}
	\\When solving a system of linear equations $Ax = b$, both Gaussian Elimination (GE) and Gauss-Jordan (GJ) require a full computational complexity of $\mathcal{O}(n^3)$ operations. If the coefficient/feature matrix $A$ stays constant but the Right-Hand Side vector (RHS, $b$) changes, both GE and GJ force you to re-run the entire algorithm from scratch, costing an additional $\mathcal{O}(n^3)$ operations for every single new vector.However, LU Decomposition (LUD) detaches the matrix $A$ from the vector $b$. It processes matrix $A$ into lower ($L$) and upper ($U$) triangular forms once at an initial cost of $\mathcal{O}(\frac{2}{3}n^3)$. After this one-time step, any changing RHS vector ($b_{new}$) can be resolved using fast substitution methods costing only $\mathcal{O}(n^2)$ operations.}
\problem Solve the system using LUD:
$$\begin{aligned}
	2x_1 + x_2 + x_3 &= 4 \\
	4x_1 + 3x_2 + 3x_3 &= 10 \\
	8x_1 + 7x_2 + 9x_3 &= 26
\end{aligned}$$
\solution
$$\begin{bmatrix} 2 & 1 & 1 \\ 4 & 3 & 3 \\ 8 & 7 & 9 \end{bmatrix} \begin{bmatrix} x_1 \\ x_2 \\ x_3 \end{bmatrix} = \begin{bmatrix} 4 \\ 10 \\ 26 \end{bmatrix}$$
Apply Gaussian elimination to reduce matrix $A$ to an upper triangular matrix $U$, while carefully recording the multipliers to build our unit lower triangular matrix $L$.\\
$R_2 \leftarrow R_2 - 2R_1 \implies \text{ The multiplier is }L_{21}=2$\\
$R_3 \leftarrow R_3 - 4R_1\implies$ The multiplier is $l_{31} = 4$.\\
$R_3 \leftarrow R_3 - 3R_2\implies$. The multiplier is $l_{32} = 3$.
$$U = \begin{bmatrix} 2 & 1 & 1 \\ 0 & 1 & 1 \\ 0 & 0 & 2 \end{bmatrix}$$
$$L = \begin{bmatrix} 1 & 0 & 0 \\ 2 & 1 & 0 \\ 4 & 3 & 1 \end{bmatrix}$$
Since $A x = b \implies L(Ux) = b$, let $Ux=y \implies Ly=b$, now solve $Ly = b$:  
$$y = \begin{bmatrix} 4,\,\, 2, \,\, 4 \end{bmatrix}$$
Now use calculated $y$ vector to solve the equation $Ux = y$
$$x = \begin{bmatrix} 1, \,\, 0, \,\, 2 \end{bmatrix}$$
\note Once the LU decomposition of the coefficient matrix (A) has been computed, it can be reused for any new right-hand-side vector. For example, if the observed target values change to $(b_{\text{new}}=\begin{bmatrix}7, 17, 43\end{bmatrix})$, there is no need to perform Gaussian elimination again.

\problem A real estate analytics firm models the sale price of a house ($y$, in lakhs of INR) as a linear function of its Size ($(x_1)$, in thousands of sq. ft.), Number of Bedrooms ($(x_2)$), and Distance from the City Center ($(x_3)$, in km):

\[
y=x_1w_1+x_2w_2+x_3w_3,
\]

where $w_1, w_2, \text{and} w_3$ are the unknown model weights.

The firm collects data from three houses:

\begin{table}[h]
	\centering
	\begin{tabular}{|l|c|c|c|}
		\hline
		\textbf{Feature} & \textbf{House 1} & \textbf{House 2} & \textbf{House 3} \\
		\hline
		Size ($x_1$, in thousands of sq. ft.) & 2 & 4 & 8 \\
		\hline
		Bedrooms ($x_2$) & 1 & 3 & 7 \\
		\hline
		Distance ($x_3$, in km) & 1 & 3 & 9 \\
		\hline
		Sale Price ($y$, in lakhs) & 5 & 11 & 21 \\
		\hline
	\end{tabular}
\end{table}

\begin{itemize}[itemsep=0.5pt, topsep=1pt]
	\item Form the corresponding system of linear equations and express it in the matrix form $(Ax=b)$.
	\item Solve the system using LU decomposition.
	\item Later, due to changes in market conditions, only the sale prices are updated while the house features remain unchanged. The new target vector is $b_{\text{new}}=
	\begin{bmatrix}
		6,\, 12,\, 26
	\end{bmatrix}.$ And hence solve the updated system $(Ax_{\text{new}}=b_{\text{new}})$.		
\end{itemize}
\solution
$$\begin{aligned}
	2w_1 + 1w_2 + 1w_3 &= 5 \\
	4w_1 + 3w_2 + 3w_3 &= 11 \\
	8w_1 + 7w_2 + 9w_3 &= 21
\end{aligned}$$
$$\begin{bmatrix} 2 & 1 & 1 \\ 4 & 3 & 3 \\ 8 & 7 & 9 \end{bmatrix} \begin{bmatrix} w_1 \\ w_2 \\ w_3 \end{bmatrix} = \begin{bmatrix} 5 \\ 11 \\ 21 \end{bmatrix}$$
$$Ax=b$$
$$U = \begin{bmatrix} 2 & 1 & 1 \\ 0 & 1 & 1 \\ 0 & 0 & 2 \end{bmatrix}, L = \begin{bmatrix} 1 & 0 & 0 \\ 2 & 1 & 0 \\ 4 & 3 & 1 \end{bmatrix}$$
Solving $Ly=b: \begin{bmatrix} 1 & 0 & 0 \\ 2 & 1 & 0 \\ 4 & 3 & 1 \end{bmatrix} \begin{bmatrix} y_1 \\ y_2 \\ y_3 \end{bmatrix} = \begin{bmatrix} 5 \\ 11 \\ 21 \end{bmatrix}\implies y = \begin{bmatrix} 5, \,\, 1, \,\, -2 \end{bmatrix}$
\\Solving $Ux = y: \begin{bmatrix} 2 & 1 & 1 \\ 0 & 1 & 1 \\ 0 & 0 & 2 \end{bmatrix} \begin{bmatrix} w_1 \\ w_2 \\ w_3 \end{bmatrix} = \begin{bmatrix} 5 \\ 1 \\ -2 \end{bmatrix}$
$x = \begin{bmatrix} w_1, \,\, w_2, \,\, w_3 \end{bmatrix} = \begin{bmatrix} 2, \,\, 2, \,\, -1 \end{bmatrix}$
\\So the linear model for 
$$y=2x_1+2x_2-x_3$$
Now for the updated target value $b_{\text{new}}$, we need to recalculate the updated model weights vector $x_{\text{new}}$ using our existing LU factorization.
\\Solving $Ly_{\text{new}} = b_{\text{new}}:\begin{bmatrix} 1 & 0 & 0 \\ 2 & 1 & 0 \\ 4 & 3 & 1 \end{bmatrix} \begin{bmatrix} y_1 \\ y_2 \\ y_3 \end{bmatrix} = \begin{bmatrix} 6 \\ 12 \\ 26 \end{bmatrix}\implies y_{\text{new}} = \begin{bmatrix} 6, \,\, 0, \,\, 2 \end{bmatrix}$\\
Solving $Ux_{\text{new}} = y_{\text{new}}: \begin{bmatrix} 2 & 1 & 1 \\ 0 & 1 & 1 \\ 0 & 0 & 2 \end{bmatrix} \begin{bmatrix} w_1 \\ w_2 \\ w_3 \end{bmatrix} = \begin{bmatrix} 6 \\ 0 \\ 2 \end{bmatrix}\implies $$x_{\text{new}} = \begin{bmatrix} w_1, \,\, w_2, \,\, w_3 \end{bmatrix} = \begin{bmatrix} 3, \,\, -1, \,\, 1 \end{bmatrix}$
\note Notice we bypassed the computationally heavy row reduction phase entirely and leap directly to basic linear substitution workflows.

\problem An e-commerce logistics firm models the total processing and shipping time of an order ($y$, in hours) as a linear function of the total Package Weight ($x_1$, in kg), Number of Fragile Items ($x_2$), and Distance to the Delivery Hub ($x_3$, in tens of km):$$\quad y = x_1w_1 + x_2w_2 + x_3w_3$$where $w_1, w_2, \text{and } w_3$ are the unknown model weights representing resource costs per unit.The firm collects performance metrics from three historical shipments:$$\begin{array}{|l|c|c|c|}
	\hline
	\textbf{Feature} & \textbf{Shipment 1} & \textbf{Shipment 2} & \textbf{Shipment 3} \\
	\hline
	\text{Weight } (x_1, \text{ in kg}) & 1 & 2 & 1 \\
	\hline
	\text{Fragile Items } (x_2) & 1 & 4 & 7 \\
	\hline
	\text{Hub Distance } (x_3, \text{ in tens of km}) & 2 & 5 & 8 \\
	\hline
	\text{Total Shipping Time } (y, \text{ in hours}) & 13 & 34 & 49 \\
	\hline
\end{array}$$Form the corresponding system of linear equations and express it in the matrix form ($Ax=b$). Solve the system using LU decomposition. Later, due to upgrades in warehouse sorting equipment, the total times change while the shipment features remain identical. The new target vector is $b_{\text{new}} = \begin{bmatrix} 9 & 28 & 45 \end{bmatrix}^T$. Solve the updated system ($Ax_{\text{new}}=b_{\text{new}}$).
\solution \textbf{Key Answers:} $$L = \begin{bmatrix} 1 & 0 & 0 \\ 2 & 1 & 0 \\ 1 & 3 & 1 \end{bmatrix}, \quad U = \begin{bmatrix} 1 & 1 & 2 \\ 0 & 2 & 1 \\ 0 & 0 & 3 \end{bmatrix}$$
$$x = \begin{bmatrix} 3 \\ 2 \\ 4 \end{bmatrix}, x_{\text{new}} = \begin{bmatrix} 1 \\ 4 \\ 2 \end{bmatrix}$$

\subsection{Linear Combinations: }
A linear combination of a set of vectors $\mathbf{v_1}, \mathbf{v_2}, \dots, \mathbf{v_n}$ is an expression formed by multiplying each vector by a scalar constant ($c_1, c_2, \dots, c_n$) and adding the results together:$$\mathbf{y} = c_1\mathbf{v_1} + c_2\mathbf{v_2} + \dots + c_n\mathbf{v_n} = \sum_{i=1}^{n} c_i\mathbf{v_i}$$
\example Consider two vectors $\mathbf{v_1}$ and $\mathbf{v_2}$ in the three-dimensional vector space ${R}^3$:$$\mathbf{v_1} = \begin{bmatrix} 1 \\ 2 \\ 3 \end{bmatrix}, \quad \mathbf{v_2} = \begin{bmatrix} 2 \\ -1 \\ 0 \end{bmatrix}$$ Let the chosen scalar multipliers be:$$c_1 = 3, \quad c_2 = -2$$ The resulting vector $\mathbf{y}$ from the linear combination:$$\mathbf{y} = c_1\mathbf{v_1} + c_2\mathbf{v_2}$$
is given by $\mathbf{y} = \begin{bmatrix} -1 \\ 8 \\ 9 \end{bmatrix}$ is a linear combination of $\mathbf{v_1}$ and $\mathbf{v_2}$.


\example Define a baseline "Average Face" ($\bar{\mathbf{v}}$) as a $100 \times 100$ pixel grid containing a smooth, centered circular glow. (See (a)).
\begin{figure}[H]
	\centering
	\includegraphics[width=0.22\textwidth]{LinCombV}\hfill
	\includegraphics[width=0.22\textwidth]{LinCombV1}\hfill
	\includegraphics[width=0.22\textwidth]{LinCombV2}\hfill
	\includegraphics[width=0.22\textwidth]{LinCombRes}
	
	\vspace{2mm}
	
	\makebox[0.22\textwidth]{(a)}
	\hfill
	\makebox[0.22\textwidth]{(b)}
	\hfill
	\makebox[0.22\textwidth]{(c)}
	\hfill
	\makebox[0.22\textwidth]{(d)}
	
	\caption{(a) Vector $\mathbf{v}$, (b) Vector $\mathbf{v}_1$, (c) Vector $\mathbf{v}_2$, and (d) Resultant vector.}
	\label{fig:lincomb_all}
\end{figure}

Flattened row-by-row, this grid becomes a 10,000-dimensional column vector ($\mathbf{v} \in {R}^{10000}$). We then establish two independent basis vectors: $\mathbf{v}_1$, a smooth left-to-right vertical shading gradient, and $\mathbf{v}_2$, a horizontal banded pattern driven by an accelerating cosine wave.By calculating the linear combination: (See (b),(c))
$$\mathbf{v}_{\text{resultant}} = \bar{\mathbf{v}} + (0.60 \cdot \mathbf{v}_1) + (-0.50 \cdot \mathbf{v}_2)$$we observe the direct geometric effect of vector algebra on image data. Applying scalar weights ($c_1 = 0.60, c_2 = -0.50$) modifies the pixels element-by-element. The final output is a clearly warped, elongated, and shifted pattern pushed toward the bottom-left of the frame, visually demonstrating how basis vector scaling reconstructs unique data points.
\example
Transforming an RGB color image to grayscale is an another classic application of a linear combination. Each pixel's new grayscale value is a unique weighted sum of its original Red, Green, and Blue channels. \\
Now create a $90 \times 90$ three-stripe image using the loops and fill first strip with red color and second with green and third with blue.
\begin{figure}[H]
	\centering
	\includegraphics[width=0.22\textwidth]{RGB1}\hfill
	\includegraphics[width=0.22\textwidth]{RGB2}\hfill
	\includegraphics[width=0.22\textwidth]{RGB3}\hfill
	\includegraphics[width=0.22\textwidth]{RGB4}
	
	\vspace{2mm}
	
	\makebox[0.22\textwidth]{(a)}
	\hfill
	\makebox[0.22\textwidth]{(b)}
	\hfill
	\makebox[0.22\textwidth]{(c)}
	\hfill
	\makebox[0.22\textwidth]{(d)}
	
	\caption{(a) RGB, (b) $0.1R+0.5G+0.9B$, (c) $0.5R+0.5G+0.5B$, and (d) $0.7R+0.1G+0.3B$.}
	\label{fig:lincomb_RGB}
\end{figure}

\problem Check whether the vector $v$ is a linear combination of the vectors $v_i$
\\
1) $v=(4, -4, 6); v_1=(1, 2, -3), v_2=(2, -4, 6), v_3=(-1, 2, -3)$
\solution Let $v=c_1 v_1+c_2 v_2+c_3 v_3$ i.e.,
$$\begin{bmatrix} 4 \\ -4 \\ 6 \end{bmatrix}=c_1 \begin{bmatrix} 1 \\ 2 \\ -3 \end{bmatrix}+c_2 \begin{bmatrix} 2 \\ -4 \\ 6 \end{bmatrix} +c_3 \begin{bmatrix} -1 \\ 2 \\ -3 \end{bmatrix}$$
$$c_1(1) + c_2(2) + c_3(-1) = 4$$$$c_1(2) + c_2(-4) + c_3(2) = -4$$$$c_1(-3) + c_2(6) + c_3(-3) = 6$$
Augmented Matrix  $[A | b]$:
$$\begin{pmatrix} 
	1 & 2 & -1 & \bigm| & 4 \\ 
	2 & -4 & 2 & \bigm| & -4 \\ 
	-3 & 6 & -3 & \bigm| & 6 
\end{pmatrix}$$
Reduce the above matrix to Echelon form:
$$\begin{pmatrix} 
	1 & 2 & -1 & \bigm| & 4 \\ 
	0 & 2 & -1 & \bigm| & 3 \\ 
	0 & 0 & 0 & \bigm| & 0 
\end{pmatrix}$$
Hence $\rho(A)=\rho(A|B)\neq \text{number of unknowns}$ which means that the above system has many solutions: the vector $v$ can be expressed as a linear combination of given set of vectors $v_i's$.
\note To find the specific combination we need to solve the above system.
\\
2) $V = \begin{pmatrix} 3 & 5 \\ 1 & -2 \end{pmatrix}; V_1 = \begin{pmatrix} 1 & 0 \\ 1 & 0 \end{pmatrix}, \quad V_2 = \begin{pmatrix} 1 & 1 \\ 0 & 1 \end{pmatrix}, \quad V_3 = \begin{pmatrix} 0 & 1 \\ -1 & -2 \end{pmatrix}, \quad V_4 = \begin{pmatrix} 0 & 0 \\ 1 & 0 \end{pmatrix}$
\solution $$A|B = \begin{pmatrix} 
	1 & 1 & 0 & 0 & \bigm| & 3 \\ 
	0 & 1 & 1 & 0 & \bigm| & 5 \\ 
	0 & 0 & -3 & 0 & \bigm| & -7 \\
	0 & 0 & 0 & 1 & \bigm| & 3 
\end{pmatrix}$$
$\rho(A)=\rho(A|B)= \text{number of unknowns}$: which means that the above system has unique solution, the vector $v$ can be expressed as a linear combination of given set of vectors $v_i's$.

\example Consider a sentence \texttt{"The quick brown fox jumps over the lazy dog. Suddenly, the dog barks loudly at a small squirrel in the yard. Many other birds chirping in the tall trees fly away in fear, startled by the sudden noisy disruption."}
\\
Now consider a sentence \texttt{The dog barks loudly}. $$\mathbf{s} = c_1\mathbf{v}_{\text{the}} + c_2\mathbf{v}_{\text{dog}} + c_3\mathbf{v}_{\text{barks}} + c_4\mathbf{v}_{\text{loudly}}$$ Let us compares two ways of choosing the scalars ($c_1, c_2, c_3, c_4$).
Case1:  The scalars are all set to $0.25$:\\
$$\mathbf{s} = 0.25\mathbf{v}_{\text{the}} + 0.25\mathbf{v}_{\text{dog}} + 0.25\mathbf{v}_{\text{barks}} + 0.25\mathbf{v}_{\text{loudly}}$$
The word "the" doesn't tell us anything about the topic, yet it contributes exactly as much to the final vector as "dog" or "barks". Hence the meaning gets diluted by filler words.
Case2: SIF Weighting:
SIF (Smooth Inverse Frequency) is an algorithm that changes the scalars based on how rare a word is. For any given word $w$, its scalar weight $c_w$ is calculated as:$$c_w = \frac{a}{a + p(w)}$$
where $p(w)$ is the estimated probability of the word appearing in the language (its total count divided by the total number of words in the dataset). $a$ is a tiny threshold parameter (typically a constant between $10^{-3}$ and $10^{-5}$, like $0.0001$). This parameter acts as a "buffer" to prevent ultra-rare words from blowing up to infinity and completely overpowering the sentence.

\begin{table}[H]
	\centering
	\caption{SIF Weights for the Toy Paragraph ($a = 0.001$, $N = 39$)}
	\label{tab:sif_weights}
	\begin{tabular}{|l|c|c|c|}
		\hline
		\textbf{Word ($w$)} & \textbf{Paragraph Count} & \textbf{Paragraph $p(w)$ (\%)} & \textbf{SIF Weight ($c_w$)} \\ \hline
		\textbf{the}        & 6                        & 15.38\%                       & 0.0065                      \\ \hline
		\textbf{dog}        & 2                        & 5.13\%                        & 0.0191                      \\ \hline
		\textbf{barks}      & 1                        & 2.56\%                        & 0.0375                      \\ \hline
		\textbf{loudly}     & 1                        & 2.56\%                        & 0.0375                      \\ \hline
	\end{tabular}
\end{table}
Now let us map the word vectors $\mathbf{v}_{\text{the}}, \mathbf{v}_{\text{dog}}, \mathbf{v}_{\text{barks}}, \mathbf{v}_{\text{loudly}}$ to another 3 dimensional vector space (toy vector space) $[\text{animal-ness}, \text{action-ness}, \text{intensity}]$.
$$v_{\text{the}}=(0.01, 0.00, 0.01)$$
$$v_{\text{dog}}=(0.9, 0.10, 0.05)$$
$$v_{\text{barks}}=(0.2, 0.85, 0.3)$$
$$v_{\text{loudly}}=(0.05, 0.2, 0.95)$$
Now take the linear combination of these vectors with the weights: we get
$$\mathbf{s}_1 = (0.4330, \,\, 0.3450, \,\, 0.3480)$$
Now consider another sentence \texttt{A canine makes a noisy disturbance} and if we do the same procedure we get
$$\mathbf{s}_2 = (0.4311, \,\, 0.3390, \,\, 0.3867)$$
Because the important keywords map to similar conceptual regions and the SIF coefficients reduce the influence of common words such as the, a, and makes, the final vectors become very similar. Therefore, the computer identifies both sentences as having the same core meaning. 

\takeaway{\textit{When you type "A canine makes a noisy disturbance" into a modern AI or search engine, it calculates that coordinate, looks at its immediate geographic neighbors in its database, finds "The dog barks loudly", and serves it to you as a perfect conceptual match}}
\textbf{Fact: }\textit{All dogs are canines, but not all canines are dogs. }

\subsection{Linear Dependency: }
A set of vectors is linearly independent if no vector in the set can be built from a combination of the others.\\
\deff A set of vectors $\{\mathbf{v}_1, \mathbf{v}_2, \dots, \mathbf{v}_n\}$ is said to be linearly independent if the only solution to the vector equation:$$c_1\mathbf{v}_1 + c_2\mathbf{v}_2 + \dots + c_n\mathbf{v}_n = \mathbf{0}$$is the trivial solution where all $c_i's$ are exactly zero:$$c_1 = 0, \quad c_2 = 0, \quad \dots, \quad c_n = 0$$
\note Consider $n$ vectors $v_1,v_2\cdot ,v_n\in R^n$, consider the linear combination:
$$c_1 \begin{pmatrix} v_{11} \\ v_{21} \\ \vdots \\ v_{m1} \end{pmatrix} + c_2 \begin{pmatrix} v_{12} \\ v_{22} \\ \vdots \\ v_{m2} \end{pmatrix} + \dots + c_n \begin{pmatrix} v_{1n} \\ v_{2n} \\ \vdots \\ v_{mn} \end{pmatrix} = \begin{pmatrix} 0 \\ 0 \\ \vdots \\ 0 \end{pmatrix}$$
Which gives
$$\begin{aligned}
	v_{11}c_1 + v_{12}c_2 + \dots + v_{1n}c_n &= 0 \\
	v_{21}c_1 + v_{22}c_2 + \dots + v_{2n}c_n &= 0 \\
	&\vdots \\
	v_{m1}c_1 + v_{m2}c_2 + \dots + v_{mn}c_n &= 0 \\
\end{aligned}$$
In matrix form:
$$\begin{pmatrix} 
	v_{11} & v_{12} & \dots & v_{1n} \\ 
	v_{21} & v_{22} & \dots & v_{2n} \\ 
	\vdots & \vdots & \ddots & \vdots \\ 
	v_{m1} & v_{m2} & \dots & v_{mn} 
\end{pmatrix}
\begin{pmatrix} c_1 \\ c_2 \\ \vdots \\ c_n \end{pmatrix} = 
\begin{pmatrix} 0 \\ 0 \\ \vdots \\ 0 \end{pmatrix}$$
When the matrix is square ($n = m$), we use the determinant to evaluate this exact system:
\begin{itemize}
	\item If $\det(A) \neq 0$: The system has only one unique solution, forcing $c_1 = c_2 = \dots = c_n = 0$. (Linearly Independent).
	\item If $\det(A) = 0$: The system has infinitely many solutions, meaning non-zero choices for $c_i$ exist that can satisfy the equations. (Linearly Dependent).
\end{itemize}

\begin{table}[H]
	\centering
	\begin{tabular}{|l|l|}
		\hline
		\textbf{Matrix Shape} & \textbf{Method}\\ \hline
		$n=m$ & Determinant\\ \hline
		$n<m$ & Row reduction/Rank\\ \hline
		$n>m$ & Linearly dependent\\ \hline
	\end{tabular}
\end{table}

\example In language, this happens when every word adds a completely fresh, unique piece of information to the sentence.\\
Let's look at three words from our first sentence:$$\mathbf{v}_{\text{dog}} = (0.90, 0.10, 0.05)$$
$$\mathbf{v}_{\text{barks}} = (0.20, 0.85, 0.30)$$
$$\mathbf{v}_{\text{loudly}} = (0.05, 0.20, 0.95)$$
$$A = \begin{pmatrix} 
	\mathbf{v}_{\text{dog}} & \mathbf{v}_{\text{barks}} & \mathbf{v}_{\text{loudly}}
\end{pmatrix} = 
\begin{pmatrix} 
	0.90 & 0.20 & 0.05 \\ 
	0.10 & 0.85 & 0.20 \\ 
	0.05 & 0.30 & 0.95 
\end{pmatrix}$$
$\det(A) = 0.655 \neq 0$ 
The matrix $A$ is invertible (non-singular).The only solution to the system is the trivial solution: $c_1 = 0, c_2 = 0, c_3 = 0$. Therefore, the vectors $\mathbf{v}_{\text{dog}}$, $\mathbf{v}_{\text{barks}}$, and $\mathbf{v}_{\text{loudly}}$ are strictly Linearly Independent.
\\
Let's look at this modified sentence: "\texttt{The canine dog barks.}"
$$\mathbf{v}_{\text{dog}} = (0.90, 0.10, 0.05)$$
$$\mathbf{v}_{\text{barks}} = (0.20, 0.85, 0.30)$$
$$\mathbf{v}_{\text{canine}} = (0.88, 0.05, 0.05)$$
$$B = \begin{pmatrix} 
	\mathbf{v}_{\text{canine}} & \mathbf{v}_{\text{dog}} & \mathbf{v}_{\text{barks}}
\end{pmatrix} = 
\begin{pmatrix} 
	0.88 & 0.90 & 0.20 \\ 
	0.05 & 0.10 & 0.85 \\ 
	0.05 & 0.05 & 0.30 
\end{pmatrix}$$
$\det(B) = 0.01325\approx 0$ Because this number is technically not exactly zero, these vectors are technically, strictly speaking, still linearly independent in a pure algebraic sense. But 0.013 is very close to zero so AI recognizes this as mathematical redundant.
\takeaway{When you ask an AI to rephrase or clean up a sentence, it is programmed to maximize information density while minimizing redundancy.
\begin{figure}[H]
	\centering
	\includegraphics[width=0.7\linewidth]{"Linear Dependent1"}
	\caption{AI removing the redundant vector}
	\label{fig:linear-dependent1}
\end{figure}
}
\subsection{Vector Space: } A vector space (or linear space) is a collection of objects called vectors, along with a set of scalars (usually real numbers), completely governed by two primary operations:\\

\begin{tabular}{cc}
	1 & Vector Addition: Combining two vectors together (${u} + {v}$). \\
	2 & Scalar Multiplication: Scaling a vector by a number ($c{v}$).
\end{tabular}
\,\\\,
\par
\textbf{Axioms of Vector Spaces}
\begin{itemize}
	\item Closure under Addition: If ${u}$ and ${v}$ are in the space, then ${u} + {v}$ must also be in the space. (You cannot add two vectors and magically land outside the space).
	\item Closure under Scalar Multiplication: If ${u}$ is in the space and $c$ is a scalar, then $c{u}$ must also be in the space.
	\item Commutativity: ${u} + {v} = {v} + {u}$
	\item Associativity: $({u} + {v}) + {w} = {u} + ({v} + {w})$
	\item Identity Element of Addition: There exists a unique zero vector ${0}$ such that ${u} + {0} = {u}$ for all vectors.
	\item Inverse Element of Addition: For every vector ${u}$, there exists an opposite vector $-{u}$ such that ${u} + (-{u}) = {0}$.
	\item Distributivity over Vector Addition: $c(\mathbf{u} + \mathbf{v}) = c\mathbf{u} + c\mathbf{v}$
	\item Distributivity over Scalar Addition: $(c + d)\mathbf{u} = c\mathbf{u} + d\mathbf{u}$
	\item Associativity of Scalar Multiplication: $c(d\mathbf{u}) = (cd)\mathbf{u}$
	\item Identity Element of Scalar Multiplication: The scalar number $1$ acts as an identity, where $1\mathbf{u} = \mathbf{u}$.
\end{itemize}

\example \textbf{1}
\begin{itemize}
	\item Set of real numbers $R$ is a vector space over the field of scalars $R$
	\item Set of $R^2$ is also a vector space over the field $R$
	\item In general set of $R^n$ is a vector space over the field $R$
	\item Set of Complex numbers $C$ is a vector space over field $R$ or $C$
	\item Set of $M_{2\times2}$ matrices of dimension $2\times2$ is a vector space over the field of scalars $R$
	\item Set of polynomials of degree $\le 2$ ($P_2$) is a vector space over the field $R$
	\item $V=\{x\in R\,|\,x>0\}$ is not a vector space over field $R$ (No Scalar Closure, No Additive Inverse)
\end{itemize}
\example \textbf{2}
\begin{itemize}
	\item Imagine a computer trying to read books. Its vocabulary consists of 10,000 unique English words (like apple, king, queen, run, sad, etc.) forms a vector space. The word "King" is represented as a vector with a 1 at the word's position and 0 everywhere else:$$\text{"King"} = [0, 0, 0, \dots, 1, \dots, 0, 0]$$
	\item In a standard grayscale image that is 512 pixels wide and 512 pixels tall, there are $512 \times 512 = 262,144$ total pixels. Each pixel holds a number representing its brightness (usually from 0 for pure black to 255 for pure white). We can unroll this entire grid of pixels into a single, massive list of 262,144 numbers. Every possible grayscale image of that exact resolution is just a different point (a unique set of coordinates) inside this same vector space.
\end{itemize}

\example If our space is ${R}^2$, we are navigating a flat, two-dimensional plane, like the map of \textit{Vector City}. To build any destination, we have to establish our coordinate system with two fundamental commands. We call this set $\mathcal{E} = \{ {e}_1, {e}_2 \}$.
\begin{itemize}
	\item Command 1 (${e}_1$): Move exactly 1 meter East (this is our X-axis). Mathematically, it is the vector $[1, 0]^T$.
	\item Command 2 (${e}_2$): Move exactly 1 meter North (this is our Y-axis). Mathematically, it is the vector $[0, 1]^T$.
\end{itemize}
\begin{figure}[H]
	\centering
	\includegraphics[width=0.5\linewidth]{GoogleMapBasis}
	\caption{Conceptual satellite road map layout demonstrating perfect rectangular Cartesian coordinate road system}
	\label{fig:googlemapbasis}
\end{figure}

\takeaway{These two navigation commands are the bases (plural of basis) for our city map: 
\begin{itemize}
	\item By scaling these two commands, we can navigate to any intersection or point in the entire city. By scaling these two commands, we can navigate to any intersection or point in the entire city. The ability to generate any vector in $R^2$ using these two vectors ($e_1$ and $e_2$) is called spanning, and the collection $\{e_1, e_2\}$ is known as the spanning set.
	\item No matter how many times ($n$ times) we execute Command 1, we will only ever move horizontally; we can never move North or South. Similarly, repeating Command 2 any number of times will only move us vertically, never horizontally East or West. Because neither command can do the job of the other, they provide entirely unique directions. This concept is called linear independence.
\end{itemize}}
\deff A set of vectors $\mathcal{B} = \{\mathbf{v}_1, \mathbf{v}_2, \dots, \mathbf{v}_n\}$ in a vector space $V$ is called a basis for $V$ if it satisfies the following two conditions:
\begin{enumerate}
	\item Linear Independence: No vector in the set can be written as a linear combination of the remaining vectors. Mathematically, the vector equation:$$c_1\mathbf{v}_1 + c_2\mathbf{v}_2 + \dots + c_n\mathbf{v}_n = \mathbf{0}$$has only the trivial solution where all scalars equal zero ($c_1 = c_2 = \dots = c_n = 0$).
	\item Spanning: The set of vectors spans the entire space, meaning every vector $\mathbf{u}$ in $V$ can be expressed as a linear combination of the vectors in $\mathcal{B}$:$$\mathbf{u} = c_1\mathbf{v}_1 + c_2\mathbf{v}_2 + \dots + c_n\mathbf{v}_n$$for some choice of scalars $c_i$.
\end{enumerate}
\example \textbf{1}
\renewcommand{\arraystretch}{1.5}
\begin{center}
	\begin{tabular}{|l|l|}
		\hline
		\textbf{Vector Space} & \textbf{Standard Basis Set ($\mathcal{B}$)}\,
		\\ 
		\hline
		\textbf{2D Real Space} $\mathbf{R}^2$ & $\left\{ \begin{bmatrix} 1 \\ 0 \end{bmatrix}, \begin{bmatrix} 0 \\ 1 \end{bmatrix} \right\}$\,
		\\ \hline
		\textbf{3D Real Space} $\mathbf{R}^3$ & $\left\{ \begin{bmatrix} 1 \\ 0 \\ 0 \end{bmatrix}, \begin{bmatrix} 0 \\ 1 \\ 0 \end{bmatrix}, \begin{bmatrix} 0 \\ 0 \\ 1 \end{bmatrix} \right\}$\,
		\\ \hline
		\textbf{$2 \times 2$ Matrices} $M_{2\times2}$ & $\left\{ \begin{bmatrix} 1 & 0 \\ 0 & 0 \end{bmatrix}, \begin{bmatrix} 0 & 1 \\ 0 & 0 \end{bmatrix}, \begin{bmatrix} 0 & 0 \\ 1 & 0 \end{bmatrix}, \begin{bmatrix} 0 & 0 \\ 0 & 1 \end{bmatrix} \right\}$\,
		\\ \hline
		\textbf{Polynomials (Degree $\le 2$)} $P_2$ & $\{x^2, x, 1\}$\, \\ \hline
	\end{tabular}
\end{center}
\renewcommand{\arraystretch}{1}
\example \textbf{2}
Imagine a vocabulary consists of 10,000 unique English words (like apple, king, queen, run, sad, etc.). To represent the word "King", the computer sets up a 10,000-dimensional vector space where every single word gets its own axis. The word "King" is represented as a vector with a 1 at the word's position and 0 everywhere else:$$\text{"King"} = [0, 0, 0, \dots, 1, \dots, 0, 0]$$
Similarly, the word "Queen" is represented by another vector with a value of 1 at the "Queen" position and 0 in all other positions. In this representation, the "King" and "Queen" axes are perpendicular (independent), implying that the two words are completely unrelated. Consequently, the computer fails to recognize that both words belong to the same semantic category of royalty. As discussed earlier, the cosine similarity between these two words should not be zero because they share a strong conceptual relationship.

To overcome this limitation, AI engineers redefine the coordinate system. Instead of assigning a separate axis to every word, they construct a smaller set of axes that represent fundamental semantic concepts. These abstract concept vectors serve as the new basis vectors, enabling words with similar meanings or contexts to be positioned closer together in the vector space.

Let us define new basis set $\{\textbf{v}_1,\textbf{v}_2,\textbf{v}_3\}$, where
\begin{itemize}
	\item Basis Vector 1 ($\mathbf{v}_1$): The concept of Royalty (High numbers mean royal, low numbers mean commoner).
	\item Basis Vector 2 ($\mathbf{v}_2$): The concept of Gender (Positive numbers mean masculine, negative numbers mean feminine).
	\item Basis Vector 3 ($\mathbf{v}_3$): The concept of Age/Power (High numbers mean adult/mature, low numbers mean young).
\end{itemize}
In our dictionary example, $\mathbf{v}_{\text{King}}$ is a data vector representing just one single word. In contrast, $\mathbf{v}_{\text{Royalty}}$ is a basis concept vector. It does not represent the single word "Royalty". Instead, $\mathbf{v}_{\text{Royalty}}$ as an evaluation checklist. It contains 10,000 numbers, where each position corresponds to a word in the dictionary, indicating how strongly that word relates to royalty. For example:

At the position for the word "King", the value is very high (e.g., 0.9).

At the position for "Queen", it is very high (e.g., 0.9).

At the position for "Crown", it is high (e.g., 0.7).

At the position for "Apple" or "Car", the value is zero (0.0).

\textbf{If we project the word "King" onto our basis, we get the coordinates:}

$\mathbf{v}_{\text{King}} \cdot \mathbf{v}_{\text{Royalty}}$ : Royalty match = 0.99

$\mathbf{v}_{\text{King}} \cdot \mathbf{v}_{\text{Gender}}$: Gender match = 0.95 (Very Masculine)

$\mathbf{v}_{\text{King}} \cdot \mathbf{v}_{\text{Age}}$: Maturity match = 0.80

Coordinates for "King": $[0.99, 0.95, 0.80]$

\textbf{If we project the word "Queen" onto our basis, we get the coordinates:}

Royalty match = 0.99

Gender match = -0.95 (Very Feminine)

Maturity match = 0.80

Coordinates for "Queen": $[0.99, -0.95, 0.80] $

\textbf{If we project the word "Boy" onto our basis}:

Coordinates for "Boy": $[0.01, 0.90, 0.10]$ (Not Royal, Masculine, Young)

Now, the word "King" is no longer a random string of text. Mathematically, it is built as a linear combination of the concept blocks:$$\text{"King"} = (0.99 \cdot \text{Royalty}) + (0.95 \cdot \text{Masculinity}) + (0.80 \cdot \text{Maturity})$$
%Because these concepts form a basis, we can do incredible vector math. If the computer takes the vector for "King", subtracts the vector for "Man" (removing masculinity), and adds the vector for "Woman" (adding femininity), it lands exactly at the coordinates for "Queen":$$\mathbf{v}_{\text{King}} - \mathbf{v}_{\text{Man}} + \mathbf{v}_{\text{Woman}} \approx \mathbf{v}_{\text{Queen}}$$
Cosine similarity of these two vector will be $0.2845$. Even after accounting for the complete opposite signs on the Gender axis, they retain a positive similarity score ($\sim 0.28$) because the shared Royalty ($0.99$) and Human ($0.80$) dimensions pull them into the same quadrant of space. In this compressed space, semantically related words share common subspace projections, allowing linear algebra operations like vector arithmetic to work:$$\mathbf{v}_{\text{King}} - \mathbf{v}_{\text{Man}} + \mathbf{v}_{\text{Woman}} \approx \mathbf{v}_{\text{Queen}}$$
In this way computer can understand both King and Queen are related to each other and they are from Royal background.

 \takeaway{There is no single "correct" basis. You can choose infinitely many different ways to build a basis for the exact same space. In fact, a basis is just a lens you choose to look through. Depending on what problem you want the AI to solve, you will pick an entirely different set of basis concepts. For example:
\begin{itemize}
	\item Option A: The "Grammar and Syntax" Basis, [Noun, Verb, Plural]
	\item Option B: The "Sentiment and Emotion" Basis,\\ $[Positivity (Happy/Joyful), Intensity (Calm vs. High Energy), Ecstatic, Sad]$
\end{itemize}
By switching to this 3D concept basis, you have stripped away 9,997 dimensions of unnecessary noise, leaving the AI with a dense, highly efficient mathematical representation of what the word "King" actually means!}
\note If a basis represents the minimum number of independent building blocks required to generate a space, then the dimension is simply the total count of those building blocks. 

In $\mathbf{R}^2$, you need 2 independent directions (East, North) to reach any location $\rightarrow$ Dimension is 2.

In $\mathbf{R}^3$, you need 3 independent directions (East, North, Up) $\rightarrow$ Dimension is 3.

For $\mathbf{M_{2\times2}}$, you need 4 independent matrix positions $\rightarrow$ Dimension is 4.


\deff Let $V$ be a finite-dimensional vector space. The dimension of $V$, denoted as $\dim(V)$, is defined as the number of vectors in any basis for $V$

\deff Let $A$ be an $m \times n$ matrix with columns $\mathbf{a}_1, \mathbf{a}_2, \dots, \mathbf{a}_n \in \mathbf{R}^m$. The Column Space of $A$, denoted $\text{Col}(A)$, is the span of its column vectors:$$\text{Col}(A) = \text{span}\{\mathbf{a}_1, \mathbf{a}_2, \dots, \mathbf{a}_n\}$$Equivalently, it is defined as:$$\text{Col}(A) = \{ \mathbf{b} \in \mathbf{R}^m \mid \mathbf{b} = A\mathbf{x} \text{ for some } \mathbf{x} \in \mathbf{R}^n \}$$

\deff A set of vectors $\{\mathbf{a}_{i_1}, \mathbf{a}_{i_2}, \dots, \mathbf{a}_{i_k}\}$ from the columns of matrix $A$ forms a basis for the column space if:
\begin{itemize}
	\item They are linearly independent.
	\item They span $\text{Col}(A)$ (every column vector in $A$ can be expressed as a linear combination of these basis vectors).
\end{itemize}

\textbf{How to Find it: } 
\begin{itemize}
	\item Reduce matrix $A$ to its Row Echelon Form (REF) using elementary row operations.
	\item Identify the pivot columns (the columns in the REF containing leading $1$s or pivots).
	\item Select the corresponding columns from the original matrix $A$ (not from the reduced form). 
	\item These original column vectors form a basis for $\text{Col}(A)$.	
\end{itemize}
\note Rank of the matrix $A$ is called the dimension of column space:$$\dim(\text{Col}(A)) = \text{Rank}(A)$$

\deff The Null Space of an $m \times n$ matrix $A$, denoted $\text{Null}(A)$, is the set of all solutions to the homogeneous equation $A\mathbf{x} = \mathbf{0}$:$$\text{Null}(A) = \{ \mathbf{x} \in \mathbf{R}^n \mid A\mathbf{x} = \mathbf{0} \}$$
\note The dimension of the null space is called the Nullity of the matrix:$$\dim(\text{Null}(A)) = \text{The number of free variables (or parameters) in the general solution}$$

\deff Let $A$ be an $m \times n$ matrix with row vectors $\mathbf{r}_1, \mathbf{r}_2, \dots, \mathbf{r}_m \in \mathbf{R}^n$. The Row Space of $A$, denoted $\text{Row}(A)$, is the span of its row vectors:$$\text{Row}(A) = \text{span}\{\mathbf{r}_1, \mathbf{r}_2, \dots, \mathbf{r}_m\}$$Equivalently, it is defined as:$$\text{Row}(A) = \{ \mathbf{y} \in \mathbf{R}^n \mid \mathbf{y} = A^T\mathbf{x} \text{ for some } \mathbf{x} \in \mathbf{R}^m \}$$

\deff A set of vectors $\{\mathbf{r}_1, \mathbf{r}_2, \dots, \mathbf{r}_k\} \subset \mathbf{R}^n$ forms a basis for the row space if:
\begin{itemize}
	\item They are linearly independent.
	\item They span $\text{Row}(A)$ (every row vector of $A$ can be expressed as a linear combination of these basis vectors).
\end{itemize}
\textbf{How to Find it:}
\begin{itemize}
	\item Reduce matrix $A$ to its Row Echelon Form (REF).
	\item The non-zero rows of the resulting echelon matrix form a linearly independent basis for $\text{Row}(A)$.
\end{itemize}
(Unlike column space, you can use the non-zero rows directly from the reduced echelon matrix because row operations preserve the row space.)

\problem Find a basis for the Column Space, Null Space and Row Space of the following $3 \times 4$ matrix $A$, and determine their dimensions ($\text{Rank}$ and $\text{Nullity}$):$$A = \begin{bmatrix} 1 & 2 & 0 & 3 \\ 2 & 4 & 1 & 8 \\ 0 & 0 & 2 & 4 \end{bmatrix}$$
\solution Reduce $A$ to Row Echelon Form :
$$A = \begin{bmatrix} 1 & 2 & 0 & 3 \\ 2 & 4 & 1 & 8 \\ 0 & 0 & 2 & 4 \end{bmatrix}\approx ~ \begin{bmatrix} \mathbf{1} & 2 & 0 & 3 \\ 0 & 0 & \mathbf{1} & 2 \\ 0 & 0 & 0 & 0 \end{bmatrix}$$
Since Column 1 and Column 3 contain the pivots in $R$, we pick Column 1 and Column 3 from the original matrix $A$.
$$\mathcal{B}_{\text{Col}(A)} = \left\{ \begin{bmatrix} 1 \\ 2 \\ 0 \end{bmatrix}, \begin{bmatrix} 0 \\ 1 \\ 2 \end{bmatrix} \right\}$$
$$\dim(\text{Col}(A)) = \text{Rank}(A) = 2 \quad \text{(since there are 2 vectors in the basis)}$$
Basis for $\text{Row}(A)$:$$\mathcal{B}_{\text{Row}(A)} = \left\{ \begin{bmatrix} 1 & 2 & 0 & 3 \end{bmatrix}, \begin{bmatrix} 0 & 0 & 1 & 2 \end{bmatrix} \right\}$$Dimension:$$\dim(\text{Row}(A)) = \text{Rank}(A) = 2$$

Let $AX=0\Rightarrow$
$$\begin{aligned} x_1 + 2x_2 + 3x_4 &= 0 \implies x_1 = -2x_2 - 3x_4 \\ x_3 + 2x_4 &= 0 \implies x_3 = -2x_4 \end{aligned}$$
$$\mathbf{x} = \begin{bmatrix} x_1 \\ x_2 \\ x_3 \\ x_4 \end{bmatrix} = \begin{bmatrix} -2x_2 - 3x_4 \\ x_2 \\ -2x_4 \\ x_4 \end{bmatrix} = x_2 \begin{bmatrix} -2 \\ 1 \\ 0 \\ 0 \end{bmatrix} + x_4 \begin{bmatrix} -3 \\ 0 \\ -2 \\ 1 \end{bmatrix}$$
Basis for $\text{Null}(A)$:$$\mathcal{B}_{\text{Null}(A)} = \left\{ \begin{bmatrix} -2 \\ 1 \\ 0 \\ 0 \end{bmatrix}, \begin{bmatrix} -3 \\ 0 \\ -2 \\ 1 \end{bmatrix} \right\}$$
$$\text{dim(Null(A))=Nullity(A)=2(since there are 2 vectors in the basis)}$$

\example Suppose our dealership lists 4 cars ($C_1, C_2, C_3, C_4$) with 5 features (rows):
\begin{itemize}
	\item Engine Capacity (Liters)
	\item Engine Capacity (CC) $\rightarrow$ (1000 $\times$ Liters)
	\item Number of Cylinders
	\item Age of Car (Years)Mileage (Thousands of Miles)
\end{itemize}
$$A = \begin{bmatrix}  2.0 & 4.0 & 1.0 & 6.0 \\  2000 & 4000 & 1000 & 6000 \\  4 & 8 & 2 & 12 \\  5 & 10 & 3 & 15 \\  50 & 100 & 20 & 150  \end{bmatrix}  \begin{matrix}  \leftarrow \text{1. Engine (L)} \\  \leftarrow \text{2. Engine (CC)} \quad (\text{Dep on Row 1}) \\  \leftarrow \text{3. Cylinders} \\  \leftarrow \text{4. Age (Years)} \\  \leftarrow \text{5. Mileage (k miles)}  \end{matrix}$$
Let's inspect the Column Relationships (Cars):

\begin{itemize}
	\item Car 1 ($C_1$): $[2.0, 2000, 4, 5, 50]^T$ $\rightarrow$ Base mid-size car.
	\item Car 2 ($C_2$): $[4.0, 4000, 8, 10, 100]^T = \mathbf{2 \times C_1}$ $\rightarrow$ Dependent! Exactly double Car 1 across all features.
	\item Car 3 ($C_3$): $[1.0, 1000, 2, 3, 20]^T$ $\rightarrow$ Independent small car.
	\item Car 4 ($C_4$): $[6.0, 6000, 12, 15, 150]^T = \mathbf{C_1 + C_2}$ (or $3 \times C_1$) $\rightarrow$ Dependent! A simple linear combination of $C_1$ and $C_2$.
\end{itemize}
Reducing matrix $A$ yields:$$\text{REF}(A) = U = \begin{bmatrix}  \mathbf{2} & 4 & 1 & 6 \\  0 & 0 & \mathbf{1} & 0 \\  0 & 0 & 0 & 0 \\  0 & 0 & 0 & 0 \\  0 & 0 & 0 & 0  \end{bmatrix}$$
The \textbf{basis of Column Space ($\text{Col}(A)$)} consists of the original pivot columns ($C_1$ and $C_3$):$$\mathcal{B}_{\text{Col}(A)} = \left\{C_1,C_3 \right\}$$
\takeaway{Although the catalog lists 4 cars, there are really only 2 fundamental car styles ($C_1$ and $C_3$).}

The \textbf{basis of Row Space ($\text{Row}(A)$)} consists of the non-zero rows of the RREF:
$$\mathcal{B}_{\text{Row}(A)} = \left\{ \begin{bmatrix} 2 & 4 & 1 & 6 \end{bmatrix}, \begin{bmatrix} 0 & 0 & 1 & 0 \end{bmatrix} \right\}$$

\takeaway{The Row Space tells you how many independent metrics are actually being measured across the inventory.
	\begin{itemize}
		\item Row 1 of REF ($\begin{bmatrix} 2 & 4 & 1 & 6 \end{bmatrix}$): Captures the engine baseline for all 4 cars. It sets Car 3's baseline multiplier to $1$
		\item Row 2 of REF ($\begin{bmatrix} 0 & 0 & 1 & 0 \end{bmatrix}$): Isolates the exact degree to which Car 3 deviates from the standard scaling pattern of the other three cars. Acts as a correction term that adjusts for Car 3's specific deviations in Age and Mileage.
	\end{itemize} 
Out of the 5 raw columns recorded in the spreadsheet (Liters, CC, Cylinders, Age, Mileage), 3 are completely redundant. You only need 2 true underlying metrics to describe every car on the lot.\\
}
To find the \textbf{Null Space ($\text{Null}(A)$)} from REF, set up $U\mathbf{x} = \mathbf{0}$ and use back-substitution:\\
Row 2:  $\implies \mathbf{x_3 = 0}$\\
Row 1: $2x_1 + 4x_2 + 1x_3 + 6x_4 = 0$, Substitute $x_3 = 0$ into Row 1: 
$\mathbf{x_1 = -2x_2 - 3x_4}$
$$\mathbf{x} = \begin{bmatrix} x_1 \\ x_2 \\ x_3 \\ x_4 \end{bmatrix} = \begin{bmatrix} -2x_2 - 3x_4 \\ x_2 \\ 0 \\ x_4 \end{bmatrix} = x_2 \begin{bmatrix} -2 \\ 1 \\ 0 \\ 0 \end{bmatrix} + x_4 \begin{bmatrix} -3 \\ 0 \\ 0 \\ 1 \end{bmatrix}$$$$\mathcal{B}_{\text{Null}(A)} = \left\{ \begin{bmatrix} -2 \\ 1 \\ 0 \\ 0 \end{bmatrix}, \begin{bmatrix} -3 \\ 0 \\ 0 \\ 1 \end{bmatrix} \right\}$$
These vectors reveal the exact linear relationships generating the redundant cars:\\
$\begin{bmatrix} -2 & 1 & 0 & 0 \end{bmatrix}^T \implies -2 C_1 + C_2 = \mathbf{0} \implies C_2 = 2C_1$\\
$\begin{bmatrix} -3 & 0 & 0 & 1 \end{bmatrix}^T \implies -3 C_1 + C_4 = \mathbf{0} \implies C_4 = 3C_1$
\takeaway{The Null Space acts as an automated "duplicate detector," showing you exactly how to express the extra cars using the primary ones.
\begin{center}
	\begin{tabular}{|l|l|}
	\hline
	\textbf{Fundamental Subspace} & \textbf{What it measures} \\
	\hline
	Column Space & Unique physical items on the lot \\
	\hline
	Row Space & Unique feature metrics \\
	\hline
	Null Space & Duplicate formulas \\
	\hline
\end{tabular}
\end{center}}

\section{Linear Transformations, Projections and Least Square Approximations:}
\subsection{Matrices Linear Transformation:} A linear transformation $T$ is a mapping (or function) between two vector spaces, $V$ and $W$, written as:$$T: V \to W$$It takes a vector $\mathbf{u}$ from $V$ as input and produces a vector $T(\mathbf{u})$ in $W$ as output, while strictly preserving the following operations:
\begin{itemize}
	\item Additivity (Preserves Addition):$$T(\mathbf{u} + \mathbf{v}) = T(\mathbf{u}) + T(\mathbf{v})$$
	\item Homogeneity (Preserves Scalar Multiplication):$$T(c\mathbf{u}) = cT(\mathbf{u})$$
\end{itemize}
\note These two rules are often combined into a single linear combination statement:$$T(c\mathbf{u} + d\mathbf{v}) = cT(\mathbf{u}) + dT(\mathbf{v})$$
Every linear transformation $T: \mathbf{R}^n \to \mathbf{R}^m$ can be represented as multiplication by an $m \times n$ standard matrix $A$:$$T(\mathbf{x}) = A\mathbf{x}$$The columns of matrix $A$ are simply the outputs of $T$ applied to the standard basis vectors ($\mathbf{e}_1, \mathbf{e}_2, \dots, \mathbf{e}_n$).

\example Consider the transformation $T: \mathbf{R}^2 \to \mathbf{R}^2$ defined by:$$T(x, y) = (2x, x + y)$$
\solution Let $\mathbf{u} = (x_1, y_1)$ and $\mathbf{v} = (x_2, y_2)$.
$$\mathbf{u} + \mathbf{v} = (x_1 + x_2, y_1 + y_2)$$
$$T(\mathbf{u} + \mathbf{v}) = (2(x_1 + x_2), (x_1 + x_2) + (y_1 + y_2)) = (2x_1 + 2x_2, x_1 + x_2 + y_1 + y_2)$$
$$T(\mathbf{u}) = (2x_1, x_1 + y_1), \quad T(\mathbf{v}) = (2x_2, x_2 + y_2)$$
$$T(\mathbf{u}) + T(\mathbf{v}) = (2x_1 + 2x_2, x_1 + y_1 + x_2 + y_2)$$
Left Side = Right Side. Additivity Holds
$$c\mathbf{u} = (cx_1, cy_1)$$
$$T(c\mathbf{u}) = (2(cx_1), cx_1 + cy_1) = (2cx_1, c(x_1 + y_1))$$
$$cT(\mathbf{u}) = c(2x_1, x_1 + y_1) = (2cx_1, c(x_1 + y_1))$$
Left Side = Right Side. Homogeneity Holds.\\
Therefore $T(x, y) = (2x, x + y)$ is a LINEAR transformation.

\example Consider the transformation $S: \mathbf{R}^2 \to \mathbf{R}^2$ defined by:$$S(x, y) = (x^2, y)$$
\solution Let $\mathbf{u} = (x_1, y_1)$ and $\mathbf{v} = (x_2, y_2)$ 
$$\mathbf{u} + \mathbf{v} = (x_1 + x_2, y_1 + y_2)$$
$$S(\mathbf{u} + \mathbf{v}) = \left( (x_1 + x_2)^2, y_1 + y_2 \right)$$
$$S(\mathbf{u}) = (x_1^2, y_1); \,\,S(\mathbf{v}) = (x_2^2, y_2)$$
$$S(\mathbf{u}) + S(\mathbf{v}) = (x_1^2 + x_2^2, \; y_1 + y_2)$$
$$S(\mathbf{u} + \mathbf{v}) \neq S(\mathbf{u}) + S(\mathbf{v})$$
Homogeneity fails generally.
Therefore $S(x, y) = (x^2, y)$ is NOT a linear transformation.

\example Find the standard matrix $A$ for the transformation $T: \mathbf{R}^3 \to \mathbf{R}^2$ defined by:$T(x, y, z) = (2x - 3y + z, x + 4z)$.
\solution Since the input vector space is $\mathbf{R}^3$, the standard basis vectors are:$$\mathbf{e}_1 = \begin{bmatrix} 1, 0, 0 \end{bmatrix}, \mathbf{e}_2 = \begin{bmatrix} 0, 1, 0 \end{bmatrix}, \mathbf{e}_3 = \begin{bmatrix} 0 , 0, 1 \end{bmatrix}$$
$$T(1, 0, 0) = \mathbf{\begin{bmatrix} 2, 1 \end{bmatrix}}$$
$$T(0, 1, 0) = \mathbf{\begin{bmatrix} -3, 0 \end{bmatrix}}$$
$$T(0, 0, 1) = \mathbf{\begin{bmatrix} 1, 4 \end{bmatrix}}$$
$$A = \begin{bmatrix} T(\mathbf{e}_1) & T(\mathbf{e}_2) & T(\mathbf{e}_3) \end{bmatrix} = \begin{bmatrix} 2 & -3 & 1 \\ 1 & 0 & 4 \end{bmatrix}$$

\textbf{Verification:}
Let $\mathbf{x} = \begin{bmatrix} x \\ y \\ z \end{bmatrix}$:
$$A\mathbf{x} = \begin{bmatrix} 2 & -3 & 1 \\ 1 & 0 & 4 \end{bmatrix} \begin{bmatrix} x \\ y \\ z \end{bmatrix} = \begin{bmatrix} 2x - 3y + z \\ x + 0y + 4z \end{bmatrix} = \begin{bmatrix} 2x - 3y + z \\ x + 4z \end{bmatrix}=T(x)$$

\example A game engine needs to convert 3D player coordinates into 2D screen coordinates using an orthographic projection transformation $T$:$$T\left( \begin{bmatrix} x \\ y \\ z \end{bmatrix} \right) = \begin{bmatrix} 2x \\ 3y \end{bmatrix}$$
\begin{figure}[H]
	\centering
	\includegraphics[width=0.4\linewidth]{TransformationGame}
	\caption{Projection from 3D coordinates to 2D coordinates}
	\label{fig:transformationgame}
\end{figure}
\begin{itemize}
	\item Find the standard transformation matrix $A$ such that $T(\mathbf{x}) = A\mathbf{x}$.
	\item Determine where an in-game character located at 3D position $\mathbf{p} = \begin{bmatrix} 5 \\ 4 \\ 100 \end{bmatrix}$ will be displayed on the 2D mini-map.
\end{itemize}
\solution The input space is $\mathbf{R}^3$, so the standard basis vectors are:$$\mathbf{e}_1 = \begin{bmatrix} 1 \\ 0 \\ 0 \end{bmatrix}, \quad \mathbf{e}_2 = \begin{bmatrix} 0 \\ 1 \\ 0 \end{bmatrix}, \quad \mathbf{e}_3 = \begin{bmatrix} 0 \\ 0 \\ 1 \end{bmatrix}$$
$$A = \begin{bmatrix} T(\mathbf{e}_1) & T(\mathbf{e}_2) & T(\mathbf{e}_3) \end{bmatrix} = \mathbf{\begin{bmatrix} 2 & 0 & 0 \\ 0 & 3 & 0 \end{bmatrix}}$$
$$T(\mathbf{p}) = \mathbf{\begin{bmatrix} 10 \\ 12 \end{bmatrix}}$$
\subsection{standard matrices for the four core linear transformations in $\mathbf{R}^2$}
\[
\begin{array}{|l|c|c|l|}
	\hline
	\mathbf{Transformation} & \mathbf{Standard\ Matrix\ } A &  \det(A) & \mathbf{Key\ Property} \\ \hline
	\text{Rotation } (\theta) & \begin{bmatrix} \cos\theta & -\sin\theta \\ \sin\theta & \cos\theta \end{bmatrix} & +1 & \text{Preserves lengths and orientation} \\ \hline
	\text{Reflection (x-axis)} & \begin{bmatrix} 1 & 0 \\ 0 & -1 \end{bmatrix} & -1 & \text{Preserves lengths, reverses orientation} \\ \hline
	\text{Reflection (y-axis)} & \begin{bmatrix} -1 & 0 \\ 0 & 1 \end{bmatrix} & -1 & \text{Preserves lengths, reverses orientation} \\ \hline
	\text{Dilation } (k > 1) & \begin{bmatrix} k & 0 \\ 0 & k \end{bmatrix} & k^2 & \text{Expands area by factor of } k^2 \\ \hline
	\text{Contraction } (0 < k < 1) & \begin{bmatrix} k & 0 \\ 0 & k \end{bmatrix} & k^2 & \text{Shrinks area by factor of } k^2 \\ \hline
\end{array}
\]
\begin{figure}[H]
	\centering
	\includegraphics[width=0.7\linewidth]{StdTransformation}
	\caption{Standard Transformations of a triangle}
	\label{fig:stdtransformation}
\end{figure}

\subsection{Coordinate Vectors and Change of Basis}
The coordinate vector of $\mathbf{x}$ relative to basis $\mathcal{B}$, denoted $[\mathbf{x}]_\mathcal{B}$, is simply the column of weights $(c_1, c_2, \dots, c_n)$ used in that linear combination:$$[\mathbf{x}]_\mathcal{B} = \begin{bmatrix} c_1 \\ c_2 \\ \vdots \\ c_n \end{bmatrix}$$
\example Suppose we have a non-standard basis $\mathcal{B} = \left\{ \begin{bmatrix} 1 \\ 1 \end{bmatrix}, \begin{bmatrix} 0 \\ 2 \end{bmatrix} \right\}$ in $\mathbf{R}^2$. If a vector in standard coordinates is $\mathbf{x} = \begin{bmatrix} 3 \\ 8 \end{bmatrix}$, what is $[\mathbf{x}]_\mathcal{B}$?
\solution Set up the vector equation:$$\begin{bmatrix} 3 \\ 8 \end{bmatrix} = c_1 \begin{bmatrix} 1 \\ 1 \end{bmatrix} + c_2 \begin{bmatrix} 0 \\ 2 \end{bmatrix}$$Solve the linear system: From Row 1: $1c_1 + 0c_2 = 3 \implies c_1 = 3$\\
From Row 2: $1c_1 + 2c_2 = 8 \implies 3 + 2c_2 = 8 \implies c_2 = 2.5$\\
The coordinate vector relative to basis $\mathcal{B}$ is:$$[\mathbf{x}]_\mathcal{B} = \begin{bmatrix} 3 \\ 2.5 \end{bmatrix}$$
\example In the above NLP example: $$\mathbf{x}_{\text{King}} = \begin{bmatrix} 1 \\ 0 \\ 0 \\ \vdots \\ 0 \end{bmatrix}$$
If we want to express $\mathbf{x}_{\text{King}}$ as a linear combination of our 3 orthonormal semantic basis vectors $\mathcal{B} = \{\mathbf{v}_{\text{Royalty}}, \mathbf{v}_{\text{Gender}}, \mathbf{v}_{\text{Age}}\}$
Each basis vector is an evaluation checklist containing 10,000 weights:
$$\mathbf{v}_{\text{Royalty}} = \begin{bmatrix} 0.99 \\ 0.99 \\ 0.70 \\ \vdots \\ 0.00 \end{bmatrix}_{\text{(King)}} , \quad \mathbf{v}_{\text{Gender}} = \begin{bmatrix} 0.95 \\ -0.95 \\ 0.00 \\ \vdots \\ 0.00 \end{bmatrix}_{\text{(King)}} , \quad \mathbf{v}_{\text{Age}} = \begin{bmatrix} 0.80 \\ 0.80 \\ 0.10 \\ \vdots \\ 0.00 \end{bmatrix}_{\text{(King)}}$$
$$\begin{bmatrix} 1 \\ 0 \\ 0 \\ \vdots \\ 0 \end{bmatrix} = c_1 \begin{bmatrix} 0.99 \\ 0.99 \\ 0.70 \\ \vdots \\ 0.00 \end{bmatrix} + c_2 \begin{bmatrix} 0.95 \\ -0.95 \\ 0.00 \\ \vdots \\ 0.00 \end{bmatrix} + c_3 \begin{bmatrix} 0.80 \\ 0.80 \\ 0.10 \\ \vdots \\ 0.00 \end{bmatrix}$$
By solving we get $$[\mathbf{x}_{\text{King}}]_\mathcal{B} = \begin{bmatrix} c_1 \\ c_2 \\ c_3 \end{bmatrix} = \begin{bmatrix} 0.99 \\ 0.95 \\ 0.80 \end{bmatrix}$$
\takeaway{By redefining our reference frame from word based axes to a set of abstract concept axes ($\mathcal{B} = \{\mathbf{v}_{\text{Royalty}}, \mathbf{v}_{\text{Gender}}, \mathbf{v}_{\text{Age}}\}$), we project high dimensional sparse vectors onto a dense, lower dimensional semantic subspace}

\subsection{Orthogonal Projections and Gram-Schmidt method}
\deff Two vectors $\mathbf{u}$ and $\mathbf{v}$ in an inner product space (such as $\mathbf{R}^n$) are orthogonal if their inner product (dot product) is zero:$$\mathbf{u} \cdot \mathbf{v} = 0$$
\begin{figure}[H]
	\centering
	\includegraphics[width=0.7\linewidth]{Orthogonality}
	\caption{Orthogonality in main Vocabulary Vector Space}
	\label{fig:orthogonality}
\end{figure}
\begin{itemize}
	\item Orthogonal Set: A set of vectors $\{\mathbf{v}_1, \mathbf{v}_2, \dots, \mathbf{v}_k\}$ is orthogonal if every pair of distinct vectors is orthogonal ($\mathbf{v}_i \cdot \mathbf{v}_j = 0$ for all $i \neq j$).
	\item Linear Independence: Any set of non-zero orthogonal vectors is automatically linearly independent.
\end{itemize}
\deff The orthogonal projection of a vector $\mathbf{y}$ onto a non-zero vector $\mathbf{u}$ decomposes $\mathbf{y}$ into two components: one component parallel to $\mathbf{u}$, and one perpendicular to $\mathbf{u}$.$$\text{proj}_{\mathbf{u}} \mathbf{y} = \left( \frac{\mathbf{y} \cdot \mathbf{u}}{\mathbf{u} \cdot \mathbf{u}} \right) \mathbf{u}$$
\deff If $W$ is a subspace spanned by an orthogonal basis $\{\mathbf{u}_1, \mathbf{u}_2, \dots, \mathbf{u}_k\}$, the orthogonal projection of $\mathbf{y}$ onto $W$ is the sum of its projections onto each basis vector:$$\text{proj}_W \mathbf{y} = \sum_{i=1}^{k} \left( \frac{\mathbf{y} \cdot \mathbf{u}_i}{\mathbf{u}_i \cdot \mathbf{u}_i} \right) \mathbf{u}_i$$
\note $\text{proj}_W \mathbf{y}$ is the closest point in $W$ to the vector $\mathbf{y}$. The error vector $(\mathbf{y} - \text{proj}_W \mathbf{y})$ is orthogonal to every vector in $W$.
\takeaway{\textbf{1. }Standard word embeddings often accidentally learn societal biases (e.g., associating certain professions with gender). Linear projections are used to remove bias without destroying semantic meaning: \\
	Identify the Gender Subspace ($W$): Compute a direction $W$ spanned by gendered difference vectors (e.g., $\mathbf{v}_{\text{he}} - \mathbf{v}_{\text{she}}$, $\mathbf{v}_{\text{man}} - \mathbf{v}_{\text{woman}}$).\\
	Project and Subtract: To debias a neutral word vector like $\mathbf{x}_{\text{doctor}}$, calculate its projection onto the gender subspace $\text{proj}_W(\mathbf{x}_{\text{doctor}})$ and subtract it:$$\mathbf{x}_{\text{debiased}} = \mathbf{x}_{\text{doctor}} - \text{proj}_W(\mathbf{x}_{\text{doctor}})$$This removes gender bias while preserving the core semantic properties of "doctor".}
	\takeaway{\textbf{2. }Dimensionality Reduction $\&$ PCA: When compressing embeddings, Principal Component Analysis (PCA - will be discussed later) identifies a lower dimensional subspace $W$ that captures the maximum variance of the text dataset. 	Projecting all document vectors onto $W$ compresses storage while retaining the vast majority of semantic relationships.}
	
	\subsection{Orthonormal bases and Gram-Schmidt method:}
\deff 	A set of vectors $\{\mathbf{q}_1, \mathbf{q}_2, \dots, \mathbf{q}_k\}$ in $\mathbf{R}^n$ is an orthonormal basis if:
	\begin{itemize}
		\item 	Orthogonal: $\mathbf{q}_i \cdot \mathbf{q}_j = \begin{cases} 1 & \text{if } i = j \\ 0 & \text{if } i \neq j \quad \end{cases}$
		\item Normalized: $\text{Unit Length: } \Vert{}\mathbf{q}_i\Vert{} = \sqrt{\mathbf{q}_i \cdot \mathbf{q}_i} = 1$
	\end{itemize}
\deff Gram-Schmidt method is an algorithm that converts any linearly independent set of vectors $\{\mathbf{x}_1, \mathbf{x}_2, \dots, \mathbf{x}_k\}$ into an orthonormal set $\{\mathbf{q}_1, \mathbf{q}_2, \dots, \mathbf{q}_k\}$ spanning the same subspace.
\\
\textbf{Procedure: }Given linearly independent vectors $\{\mathbf{x}_1, \mathbf{x}_2, \dots, \mathbf{x}_k\}$, construct the orthonormal basis $\{\mathbf{q}_1, \mathbf{q}_2, \dots, \mathbf{q}_k\}$ step-by-step:
\begin{itemize}
	\item Step 1: $\mathbf{v}_1 = \mathbf{x}_1 \implies \mathbf{q}_1 = \frac{\mathbf{v}_1}{\sqrt{\mathbf{v}_1 \cdot \mathbf{v}_1}}$
	\item Step 2: $\mathbf{v}_2 = \mathbf{x}_2 - (\mathbf{x}_2 \cdot \mathbf{q}_1)\mathbf{q}_1 \implies \mathbf{q}_2 = \frac{\mathbf{v}_2}{\sqrt{\mathbf{v}_2 \cdot \mathbf{v}_2}}$
	\item Step $k$: $\mathbf{v}_k = \mathbf{x}_k - \sum_{j=1}^{k-1} (\mathbf{x}_k \cdot \mathbf{q}_j)\mathbf{q}_j \implies \mathbf{q}_k = \frac{\mathbf{v}_k}{\sqrt{\mathbf{v}_k \cdot \mathbf{v}_k}}$
\end{itemize}

\deff QR Factorization decomposes a matrix $A$ into the product of an orthogonal (or orthonormal) matrix $Q$ and an upper triangular matrix $R$:$$A = QR$$
\begin{itemize}
	\item $Q$ (Orthonormal Columns): An $m \times n$ matrix whose columns $q_1, q_2, \dots, q_n$ are mutually perpendicular unit vectors ($Q^T Q = I$).
	\item $R$ (Upper Triangular): An $n \times n$ upper triangular matrix containing the inner products (projection coefficients) that reconstruct $A$ from $Q$.
\end{itemize}
\textbf{Procedure: } Given a matrix $A = \begin{pmatrix} a_1 & a_2 & \cdots & a_n \end{pmatrix}$ with linearly independent columns:
\begin{itemize}
	\item \textbf{Step 1:} Orthogonalize $\&$ Normalize ($Q$): Apply the Gram-Schmidt process to columns $a_j$ to obtain unit orthogonal vectors $q_j$.
	\item \textbf{Step 2:} Because the columns of $Q$ are orthonormal, the dot product of any column with itself is $1$, and with any other column is $0$. This identity property is expressed as:$$Q^T Q = I$$
	Hence, the Projection Coefficients ($R$) is given by: $R = Q^T A$ 
\end{itemize}
\example Factorize the matrix $A$ in to $QR$
$$A = \begin{pmatrix} 1 & 1 \\ 1 & 0 \\ 0 & 1 \end{pmatrix}$$
\solution $$\mathbf{x}_1 = \begin{pmatrix} 1 \\ 1 \\ 0 \end{pmatrix}, \quad \mathbf{x}_2 = \begin{pmatrix} 1 \\ 0 \\ 1 \end{pmatrix}$$Step 1$$\mathbf{v}_1 = \mathbf{x}_1 = \begin{pmatrix} 1 \\ 1 \\ 0 \end{pmatrix}\Rightarrow \mathbf{q}_1 = \frac{\mathbf{v}_1}{\sqrt{\mathbf{v}_1 \cdot \mathbf{v}_1}} = \frac{1}{\sqrt{2}} \begin{pmatrix} 1 \\ 1 \\ 0 \end{pmatrix} = \begin{pmatrix} 1/\sqrt{2} \\ 1/\sqrt{2} \\ 0 \end{pmatrix}$$
Step 2$$\mathbf{x}_2 \cdot \mathbf{q}_1 = \begin{pmatrix} 1 \\ 0 \\ 1 \end{pmatrix} \cdot \begin{pmatrix} 1/\sqrt{2} \\ 1/\sqrt{2} \\ 0 \end{pmatrix} = \frac{1}{\sqrt{2}}$$
$$\mathbf{v}_2 = \mathbf{x}_2 - (\mathbf{x}_2 \cdot \mathbf{q}_1)\mathbf{q}_1 = \begin{pmatrix} 1 \\ 0 \\ 1 \end{pmatrix} - \frac{1}{\sqrt{2}} \begin{pmatrix} 1/\sqrt{2} \\ 1/\sqrt{2} \\ 0 \end{pmatrix} = \begin{pmatrix} 1 \\ 0 \\ 1 \end{pmatrix} - \begin{pmatrix} 1/2 \\ 1/2 \\ 0 \end{pmatrix} = \begin{pmatrix} 1/2 \\ -1/2 \\ 1 \end{pmatrix}$$$$\sqrt{\mathbf{v}_2 \cdot \mathbf{v}_2} = \sqrt{\left(\frac{1}{2}\right)^2 + \left(-\frac{1}{2}\right)^2 + 1^2} = \sqrt{\frac{1}{4} + \frac{1}{4} + 1} = \sqrt{\frac{3}{2}} = \frac{\sqrt{6}}{2}$$
$$\mathbf{q}_2 = \frac{\mathbf{v}_2}{\sqrt{\mathbf{v}_2 \cdot \mathbf{v}_2}} = \frac{2}{\sqrt{6}} \begin{pmatrix} 1/2 \\ -1/2 \\ 1 \end{pmatrix} = \begin{pmatrix} 1/\sqrt{6} \\ -1/\sqrt{6} \\ 2/\sqrt{6} \end{pmatrix}$$
$$Q = \begin{pmatrix} \mathbf{q}_1 & \mathbf{q}_2 \end{pmatrix} = \begin{pmatrix} 1/\sqrt{2} & 1/\sqrt{6} \\ 1/\sqrt{2} & -1/\sqrt{6} \\ 0 & 2/\sqrt{6} \end{pmatrix}$$
$R = Q^T A$
$$R = \begin{pmatrix} 1/\sqrt{2} & 1/\sqrt{2} & 0 \\ 1/\sqrt{6} & -1/\sqrt{6} & 2/\sqrt{6} \end{pmatrix} \begin{pmatrix} 1 & 1 \\ 1 & 0 \\ 0 & 1 \end{pmatrix}$$
$$R = \begin{pmatrix} \sqrt{2} & 1/\sqrt{2} \\ 0 & \sqrt{6}/2 \end{pmatrix}$$
$$A = \begin{pmatrix} 1 & 1 \\ 1 & 0 \\ 0 & 1 \end{pmatrix}=\begin{pmatrix} 1/\sqrt{2} & 1/\sqrt{6} \\ 1/\sqrt{2} & -1/\sqrt{6} \\ 0 & 2/\sqrt{6} \end{pmatrix}\begin{pmatrix} \sqrt{2} & 1/\sqrt{2} \\ 0 & \sqrt{6}/2 \end{pmatrix}=QR$$

\deff Least squares approximation solves an overdetermined system $Ax = b$ when no exact solution exists by finding a vector $\hat{x}$ that minimizes the squared Euclidean distance $\Vert{}A\hat{x} - b\Vert{}^2$. Geometrically, it projects the target vector $b$ orthogonally onto the column space of $A$.\\
When $b$ lies outside the column space of $A$, the error vector $e = b - A\hat{x}$ cannot be zero. To minimize the length of $e$, $e$ must be orthogonal (perpendicular) to every column in $A$:\\
For $e$ to be perpendicular to every vector in $\text{Col}(A)$, it must be perpendicular to each column vector $a_i$:$$a_1 \cdot e = 0, \quad a_2 \cdot e = 0, \quad \dots, \quad a_n \cdot e = 0$$
This gives a system of equations 
$$\begin{pmatrix} a_1^T e \\ a_2^T e \\ \vdots \\ a_n^T e \end{pmatrix} = \begin{pmatrix} 0 \\ 0 \\ \vdots \\ 0 \end{pmatrix}$$
or simply
$$A^T e = 0 \implies A^T (b - A\hat{x}) = 0$$
$$A^T A \hat{x} = A^T b$$
$$\hat{x} = (A^T A)^{-1} A^T b$$
In practice, computing $(A^T A)^{-1}$ directly causes numerical instability. Substituting $A = QR$ simplifies the normal equations significantly:
$$(QR)^T (QR) \hat{x} = (QR)^T b$$
Since $Q^T Q = I$:
$$R \hat{x} = Q^T b$$
Because $R$ is upper triangular, $\hat{x}$ is solved quickly and accurately using simple back-substitution.

\example Find the least squares solution $\hat{x}$ for $Ax = b$: $$A = \begin{pmatrix} 1 & 1 \\ 1 & 0 \\ 0 & 1 \end{pmatrix}, \quad b = \begin{pmatrix} 4 \\ 1 \\ 2 \end{pmatrix}$$
\solution The $Q$ and $R$ matrices are computed in the previous example:
$$Q=\begin{pmatrix} 1/\sqrt{2} & 1/\sqrt{6} \\ 1/\sqrt{2} & -1/\sqrt{6} \\ 0 & 2/\sqrt{6} \end{pmatrix}; R=\begin{pmatrix} \sqrt{2} & 1/\sqrt{2} \\ 0 & \sqrt{6}/2 \end{pmatrix}$$

$$Q^T b = \begin{pmatrix} 1/\sqrt{2} & 1/\sqrt{2} & 0 \\ 1/\sqrt{6} & -1/\sqrt{6} & 2/\sqrt{6} \end{pmatrix} \begin{pmatrix} 4 \\ 1 \\ 2 \end{pmatrix} = \begin{pmatrix} \frac{4}{\sqrt{2}} + \frac{1}{\sqrt{2}} + 0 \\ \frac{4}{\sqrt{6}} - \frac{1}{\sqrt{6}} + \frac{4}{\sqrt{6}} \end{pmatrix} = \begin{pmatrix} 5/\sqrt{2} \\ 7/\sqrt{6} \end{pmatrix}$$
Solve $R \hat{x} = Q^T b$:$$\begin{pmatrix} \sqrt{2} & 1/\sqrt{2} \\ 0 & \sqrt{6}/2 \end{pmatrix} \begin{pmatrix} \hat{x}_1 \\ \hat{x}_2 \end{pmatrix} = \begin{pmatrix} 5/\sqrt{2} \\ 7/\sqrt{6} \end{pmatrix}$$
$$\hat{x} = \begin{pmatrix} 4/3 \\ 7/3 \end{pmatrix}$$


\example Suppose we want an AI algorithm to predict a movie's star rating ($y$, from 1 to 5) automatically by counting key sentiment words in review text. We model the star rating as a linear combination of word counts:
$$\text{Rating} = w_1 \cdot (\text{count of "great"}) + w_2 \cdot (\text{count of "boring"}) + b$$
We need to solve for three unknown parameters: word weights $w_1, w_2$ and the base rating $b$.
\begin{itemize}
	\item 1. Collect Data (The Reviews): We sample 5 customer reviews:
	\begin{itemize}
		\item Review 1: 2 "great", 0 "boring" $\to$ 5 Stars
		\item Review 2: 0 "great", 2 "boring" $\to$ 1 Star
		\item Review 3: 1 "great", 1 "boring" $\to$ 3 Stars
		\item Review 4: 2 "great", 1 "boring" $\to$ 4 Stars
		\item Review 5: 0 "great", 1 "boring" $\to$ 2 Stars
	\end{itemize}
	\item 2. Plugging these word counts into your formula yields the linear system $A\mathbf{w} \approx \mathbf{b}$:
	$$\begin{aligned} 2w_1 + 0w_2 + b &= 5.0 \\ 0w_1 + 2w_2 + b &= 1.0 \\ 1w_1 + 1w_2 + b &= 3.0 \\ 2w_1 + 1w_2 + b &= 4.0 \\ 0w_1 + 1w_2 + b &= 2.0 \end{aligned}$$
	3 Unknowns ($n=3$): Weights $w_1, w_2$ and baseline $b$. \\ 5 Equations ($m=5$): 5 real-world reviews. \\ Because $5 > 3$ (far more review equations than unknown word weights), the system is overdetermined.
	\item Human language is inherently noisy and subjective. Notice that Review 2 gave 1 star for 2 "boring"s, but Review 5 gave 2 stars for 1 "boring".These equations contradict one another. No single combination of $(w_1, w_2, b)$ can draw a plane that passes through all 5 review data points perfectly. An exact solution does not exist.
	\item Least Squares finds the single best weight vector $\hat{\mathbf{w}}$ by minimizing the total squared prediction error across all reviews:$$\min_{w_1, w_2, b} \sum_{i=1}^{5} \left( \text{Predicted Rating}_i - \text{Actual Rating}_i \right)^2$$Result: It outputs the best statistical weights (e.g., $w_1 = +1.1$, $w_2 = -0.9$, $b = 2.8$).
	\takeaway{Without human supervision telling the computer what words mean, least squares optimization assigns positive values to praise (e.g., $w_1 = +1.1$ for "great") and negative values to criticism (e.g., $w_2 = -0.9$ for "boring").}
	
	\section{Eigenvalues- Eigenvectors, Matrix Decomposition and Principal Component Analysis}
	Geometrically, a matrix acts as a spatial transformation that rotates, stretches, or squishes space. Eigenvectors are the special directions that do not rotate when transformed by a matrix, and eigenvalues are the scaling factors that tell you how much those vectors stretch, shrink, or flip.
	\deff The Fundamental Equation$$A \mathbf{v} = \lambda \mathbf{v}$$
	\begin{itemize}
		\item $A$ (Matrix): The linear transformation operator ($n \times n$).
		\item $\mathbf{v}$ (Eigenvector): A non-zero vector whose directional orientation remains invariant after transformation.
		\item $\lambda$ (Eigenvalue): A scalar multiplier representing how much $\mathbf{v}$ is scaled.
	\end{itemize}
	
	\example Given Matrix$$A = \begin{pmatrix} 2 & 1 & 1 \\ 1 & 2 & 1 \\ 1 & 1 & 2 \end{pmatrix}$$
	\solution
	$$\det(A - \lambda I) = \begin{vmatrix} 2-\lambda & 1 & 1 \\ 1 & 2-\lambda & 1 \\ 1 & 1 & 2-\lambda \end{vmatrix} = 0$$
	$\Rightarrow \lambda_1 = 4, \lambda_2 = 1, \quad \lambda_3 = 1$ are the eigenvalues
	\begin{itemize}
		\item[$\bullet$] $\mathbf{\lambda = 4:} $ \\
		$$\begin{aligned} -2x_1 + x_2 + x_3 &= 0 \\ x_1 - 2x_2 + x_3 &= 0 \\ x_1 + x_2 - 2x_3 &= 0 \end{aligned}$$
		$$\mathbf{v}_1 = \begin{pmatrix} 1 \\ 1 \\ 1 \end{pmatrix}$$
		\item[$\bullet$] $\mathbf{\lambda = 1:}$\\
		$$x_1 + x_2 + x_3 = 0 \implies x_1 = -x_2 - x_3$$
		Because you have 1 equation and 3 variables, there are 2 free variables ($x_2$ and $x_3$). This means the repeated eigenvalue generates two linearly independent eigenvectors.\\
		Set $x_2 = 1, x_3 = 0 \Rightarrow \mathbf{v}_2 = \begin{pmatrix} -1 \\ 1 \\ 0 \end{pmatrix}$. 
		Set $x_2 = 0, x_3 = 1 \Rightarrow \mathbf{v}_3 = \begin{pmatrix} -1 \\ 0 \\ 1 \end{pmatrix}$		
	\end{itemize}
	\example 
	Consider a matrix $A = \begin{pmatrix} 2 & 1 \\ 1 & 2 \end{pmatrix}$, eigenvalues are $$\mathbf{\lambda_1 = 3}, \quad \mathbf{\lambda_2 = 1}$$ and corresponding eigenvectors are 
	$ \mathbf{v}_1 = \begin{pmatrix} 1 \\ 1 \end{pmatrix},\mathbf{v}_2 = \begin{pmatrix} -1 \\ 1 \end{pmatrix} $
	The red circle represents all unit vectors (vectors with length = 1) as input. The blue ellipse represents what happens when the matrix $A$ acts on space. It distorts and stretches the red unit circle into an elongated ellipse. 	
		\begin{figure}[H]
		\centering
		\begin{minipage}{0.3\textwidth}
			\centering
			\includegraphics[width=\linewidth]{eigenValue01}
		\end{minipage}
		\hspace{0.01\textwidth}
		\begin{minipage}{0.3\textwidth}
			\centering
			\includegraphics[width=\linewidth]{eigenValue02}
		\end{minipage}
		
		\caption{Visualization of what eigenvector}
		\label{fig:eigenValues}
	\end{figure}
	When matrix $A$ transforms $v$, it becomes the green vector $Av$. Notice that $Av$ changed both its length and its direction. But if we watch closely, there are a few special directions where the vectors do not rotate at all—they just get pushed straight out or squashed along the exact same line.
	
	\subsection{Symmetric matrices:}
	A real square matrix $A$ is symmetric if it is equal to its transpose:$$A = A^T$$
	Key Properties:
	\begin{itemize}
		\item Real Eigenvalues: All eigenvalues ($\lambda_i$) of a real symmetric matrix are guaranteed to be real numbers.
		\item Orthogonal Eigenvectors: Eigenvectors corresponding to distinct eigenvalues are mutually orthogonal.
	\end{itemize}
	\deff A matrix $A$ is said to be an orthogonal matrix if  $A^T A = I$, or $A^{-1} = A^T$.
	\deff By the Spectral Theorem, any real symmetric matrix $A \in \mathbf{R}^{n \times n}$ can be orthogonally diagonalized:$$A = Q \Lambda Q^T$$
	
	where: $Q$ is an orthogonal matrix whose columns are orthonormal eigenvectors $\mathbf{q}_1, \mathbf{q}_2, \dots, \mathbf{q}_n$ (meaning $Q^T Q = I$, or $Q^{-1} = Q^T$). \\
	$\Lambda$ is a diagonal matrix containing the real eigenvalues $\lambda_1, \lambda_2, \dots, \lambda_n$ along its main diagonal.
	
	
	
	\subsection{Orthogonal Diagonalization:} Given an $n \times n$ real symmetric matrix $A$, we have to find an orthogonal matrix $Q$ and a diagonal matrix $\Lambda$ such that $A = Q \Lambda Q^T$.
	\begin{itemize}
		\item Step 1: Compute the Eigenvalues $\lambda_1, \lambda_2, \dots, \lambda_n$.
		\item Step 2: Find the Basis of Eigenvectors $v_i's$ for Each Eigenvalue $\lambda_i's$.
		\item Step 3: Orthonormalize the Eigenvectors:\\ \textbf{Case A: }Eigenvectors corresponding to different eigenvalues are automatically orthogonal. We only need to normalize each eigenvector by dividing by its magnitude:
		$$\mathbf{q}_i = \frac{\mathbf{v}_i}{\Vert{}\mathbf{v}_i\Vert{}}$$
		\textbf{Case B: }Eigenvectors from the same eigenspace (i.e., Repeated Eigenvalues) are not guaranteed to be mutually orthogonal. Apply the Gram-Schmidt process to the basis vectors of that eigenspace to make them orthogonal.Normalize all resulting orthogonal vectors to unit length:
		$$\mathbf{q}_i = \frac{\mathbf{u}_i}{\Vert{}\mathbf{u}_i\Vert{}}$$.
		\item Step 4: Construct Matrices $Q = \begin{bmatrix} \mathbf{q}_1 & \mathbf{q}_2 & \dots & \mathbf{q}_n \end{bmatrix}$ and $\Lambda = \begin{bmatrix} \lambda_1 & 0 & \dots & 0 \\ 0 & \lambda_2 & \dots & 0 \\ \vdots & \vdots & \ddots & \vdots \\ 0 & 0 & \dots & \lambda_n \end{bmatrix}$		
	\end{itemize}
	
	\example Orthogonally diagonalize $A = \begin{bmatrix} 4 & 2 \\ 2 & 1 \end{bmatrix}$ into $A = Q \Lambda Q^T$.
	\solution Step 1: Eigenvalues $\Rightarrow \det(A - \lambda I) = 0 \implies \lambda^2 - 5\lambda = 0 \implies \lambda_1 = 5, \quad \lambda_2 = 0$\\
	Step 2: Eigenvectors : $(A - \lambda I)\mathbf{v} = \mathbf{0}\implies \mathbf{v}_1 = \begin{bmatrix} 2 \\ 1 \end{bmatrix}, \mathbf{v}_2 = \begin{bmatrix} -1 \\ 2 \end{bmatrix}$\\
	Step 3: Normalize Eigenvectors : $(\mathbf{q}_i = \frac{\mathbf{v}_i}{\Vert{}\mathbf{v}_i\Vert{}}) \implies \mathbf{q}_1 = \frac{1}{\sqrt{5}}\begin{bmatrix} 2 \\ 1 \end{bmatrix}, \mathbf{q}_2 = \frac{1}{\sqrt{5}}\begin{bmatrix} -1 \\ 2 \end{bmatrix}$ 
	Step 4: $$Q = \begin{bmatrix} \frac{2}{\sqrt{5}} & -\frac{1}{\sqrt{5}} \\ \frac{1}{\sqrt{5}} & \frac{2}{\sqrt{5}} \end{bmatrix}, \quad \Lambda = \begin{bmatrix} 5 & 0 \\ 0 & 0 \end{bmatrix}$$
	\example Orthogonally diagonalize $A = \begin{bmatrix} 3 & 0 & 0 \\ 0 & 2 & 1 \\ 0 & 1 & 2 \end{bmatrix}$ into $A = Q \Lambda Q^T$
	\solution Step 1: Eigenvalues ($\det(A - \lambda I) = 0 \implies \lambda_1 = 3, 3 \text{ (multiplicity 2)}, \quad \lambda_2 = 1$\\
	Step 2: Eigenvectors ($(A - \lambda I)\mathbf{v} = \mathbf{0}$) For $\lambda_1 = 3$:$$\begin{bmatrix} 0 & 0 & 0 \\ 0 & -1 & 1 \\ 0 & 1 & -1 \end{bmatrix}\begin{bmatrix} x_1 \\ x_2 \\ x_3 \end{bmatrix} = \mathbf{0} \implies -x_2 + x_3 = 0 \implies x_2 = x_3, \; x_1 \text{ is free}$$
	Two linearly independent eigenvectors:$$\mathbf{v}_1 = \begin{bmatrix} 1 \\ 0 \\ 0 \end{bmatrix}, \quad \mathbf{v}_2 = \begin{bmatrix} 0 \\ 1 \\ 1 \end{bmatrix}$$
	For $\lambda_2 = 1$:$$\begin{bmatrix} 2 & 0 & 0 \\ 0 & 1 & 1 \\ 0 & 1 & 1 \end{bmatrix}\begin{bmatrix} x_1 \\ x_2 \\ x_3 \end{bmatrix} = \mathbf{0} \implies x_1 = 0 \text{ and } x_2 + x_3 = 0 \implies \mathbf{v}_3 = \begin{bmatrix} 0 \\ -1 \\ 1 \end{bmatrix}$$
	Step 3: Normalize Eigenvectors ($\mathbf{q}_i = \frac{\mathbf{v}_i}{\Vert{}\mathbf{v}_i\Vert{}}$) 
	$\implies \mathbf{q}_1 = \begin{bmatrix} 1 \\ 0 \\ 0 \end{bmatrix}, \mathbf{q}_2 = \frac{1}{\sqrt{2}}\begin{bmatrix} 0 \\ 1 \\ 1 \end{bmatrix}, \mathbf{q}_3 = \frac{1}{\sqrt{2}}\begin{bmatrix} 0 \\ -1 \\ 1 \end{bmatrix}$ \\
	Step 4: Final Matrices ($Q$ and $\Lambda$)$$Q = \begin{bmatrix} 1 & 0 & 0 \\ 0 & \frac{1}{\sqrt{2}} & -\frac{1}{\sqrt{2}} \\ 0 & \frac{1}{\sqrt{2}} & \frac{1}{\sqrt{2}} \end{bmatrix}, \quad \Lambda = \begin{bmatrix} 3 & 0 & 0 \\ 0 & 3 & 0 \\ 0 & 0 & 1 \end{bmatrix}$$
	
	\subsection{Vector Expansion in the Eigenvector Basis}
	Because the orthonormal eigenvectors $\{\mathbf{q}_1, \mathbf{q}_2, \dots, \mathbf{q}_n\}$ form a complete basis for $\mathbf{R}^n$, any arbitrary vector $\mathbf{x}$ can be uniquely written as a linear combination of these basis vectors:$$\mathbf{x} = c_1\mathbf{q}_1 + c_2\mathbf{q}_2 + \dots + c_n\mathbf{q}_n$$To find each coordinate coefficient $c_i$, we take the inner product (dot product) of $\mathbf{x}$ with $\mathbf{q}_i$:
	$$c_i = \mathbf{q}_i^T \mathbf{x}$$
	Thus, expanding $\mathbf{x}$ in terms of the eigenvector basis gives:$$\mathbf{x} = (\mathbf{q}_1^T \mathbf{x})\mathbf{q}_1 + (\mathbf{q}_2^T \mathbf{x})\mathbf{q}_2 + \dots + (\mathbf{q}_n^T \mathbf{x})\mathbf{q}_n = \sum_{i=1}^n (\mathbf{q}_i^T \mathbf{x})\mathbf{q}_i$$
	Now, apply the linear transformation $A$ to the vector $\mathbf{x}$:$$A\mathbf{x} = A \left( \sum_{i=1}^n (\mathbf{q}_i^T \mathbf{x})\mathbf{q}_i \right) = \sum_{i=1}^n (\mathbf{q}_i^T \mathbf{x}) A\mathbf{q}_i$$Since each $\mathbf{q}_i$ is an eigenvector of $A$ satisfying $A\mathbf{q}_i = \lambda_i \mathbf{q}_i$:$$A\mathbf{x} = \sum_{i=1}^n \lambda_i (\mathbf{q}_i^T \mathbf{x}) \mathbf{q}_i$$
	By factoring out $\mathbf{x}$ from the sum above, the matrix $A$ itself can be expanded into a sum of rank-1 projection matrices:$$A = \lambda_1 \mathbf{q}_1 \mathbf{q}_1^T + \lambda_2 \mathbf{q}_2 \mathbf{q}_2^T + \dots + \lambda_n \mathbf{q}_n \mathbf{q}_n^T = \sum_{i=1}^n \lambda_i (\mathbf{q}_i \mathbf{q}_i^T)$$
	\begin{itemize}
		\item $\mathbf{q}_i \mathbf{q}_i^T$ is an orthogonal projection operator that projects any vector onto the 1D subspace spanned by the basis vector $\mathbf{q}_i$.
		\item $\lambda_i$ acts as a scaling factor along that specific basis direction.
	\end{itemize}
	\insight{When a symmetric matrix $A$ transforms a vector $\mathbf{x}$ via $A\mathbf{x} = Q \Lambda Q^T \mathbf{x}$, it executes three distinct steps using the eigenvector basis:
		\begin{itemize}
			\item Change of Basis ($Q^T \mathbf{x}$): Projects the input vector $\mathbf{x}$ onto the new orthogonal coordinate axes (the eigenvectors $\mathbf{q}_i$). This computes the scalar coordinates $c_i = \mathbf{q}_i^T \mathbf{x}$ along each eigenvector axis.
			\item Independent Scaling ($\Lambda$): Stretches or compresses each coordinate independently by multiplying it by its corresponding eigenvalue $\lambda_i$.
			\item Reconstruction / Change Back to Standard Basis ($Q$): Recombines the scaled components along the direction of each eigenvector to yield the transformed vector in the original coordinate system.
		\end{itemize}
	}
	\example Consider $$A = \begin{bmatrix} 3 & 1 \\ 1 & 3 \end{bmatrix}$$\\
	Eigenvalue $\lambda_1 = 4$ with normalized eigenvector $\mathbf{q}_1 = \begin{bmatrix} \frac{1}{\sqrt{2}} \\ \frac{1}{\sqrt{2}} \end{bmatrix}$\\Eigenvalue $\lambda_2 = 2$ with normalized eigenvector $\mathbf{q}_2 = \begin{bmatrix} -\frac{1}{\sqrt{2}} \\ \frac{1}{\sqrt{2}} \end{bmatrix}$
	$$A = \lambda_1 (\mathbf{q}_1 \mathbf{q}_1^T) + \lambda_2 (\mathbf{q}_2 \mathbf{q}_2^T)$$
	Component 1 (Along $\mathbf{q}_1$):
	$$\mathbf{q}_1 \mathbf{q}_1^T = \begin{bmatrix} \frac{1}{\sqrt{2}} \\ \frac{1}{\sqrt{2}} \end{bmatrix} \begin{bmatrix} \frac{1}{\sqrt{2}} & \frac{1}{\sqrt{2}} \end{bmatrix} = \begin{bmatrix} \frac{1}{2} & \frac{1}{2} \\ \frac{1}{2} & \frac{1}{2} \end{bmatrix}$$
	Scaling by $\lambda_1 = 4$:$$\lambda_1 (\mathbf{q}_1 \mathbf{q}_1^T) = 4 \begin{bmatrix} \frac{1}{2} & \frac{1}{2} \\ \frac{1}{2} & \frac{1}{2} \end{bmatrix} = \begin{bmatrix} 2 & 2 \\ 2 & 2 \end{bmatrix}$$
	Component 2 (Along $\mathbf{q}_2$):$$\mathbf{q}_2 \mathbf{q}_2^T = \begin{bmatrix} -\frac{1}{\sqrt{2}} \\ \frac{1}{\sqrt{2}} \end{bmatrix} \begin{bmatrix} -\frac{1}{\sqrt{2}} & \frac{1}{\sqrt{2}} \end{bmatrix} = \begin{bmatrix} \frac{1}{2} & -\frac{1}{2} \\ -\frac{1}{2} & \frac{1}{2} \end{bmatrix}$$Scaling by $\lambda_2 = 2$:$$\lambda_2 (\mathbf{q}_2 \mathbf{q}_2^T) = 2 \begin{bmatrix} \frac{1}{2} & -\frac{1}{2} \\ -\frac{1}{2} & \frac{1}{2} \end{bmatrix} = \begin{bmatrix} 1 & -1 \\ -1 & 1 \end{bmatrix}$$
	Reconstructing $A$:$$\lambda_1 (\mathbf{q}_1 \mathbf{q}_1^T) + \lambda_2 (\mathbf{q}_2 \mathbf{q}_2^T) = \begin{bmatrix} 2 & 2 \\ 2 & 2 \end{bmatrix} + \begin{bmatrix} 1 & -1 \\ -1 & 1 \end{bmatrix} = \begin{bmatrix} 3 & 1 \\ 1 & 3 \end{bmatrix} = A$$
	Now let's see how $A$ transforms a sample input vector $\mathbf{x} = \begin{bmatrix} 5 \\ 1 \end{bmatrix}$.Direct Calculation:$$A\mathbf{x} = \begin{bmatrix} 3 & 1 \\ 1 & 3 \end{bmatrix} \begin{bmatrix} 5 \\ 1 \end{bmatrix} = \begin{bmatrix} 16 \\ 8 \end{bmatrix}$$
	Interpretation via Eigenvector Basis Expansion: \\
	Project $\mathbf{x}$ onto the Eigenvector Basis (Get Coordinates). Compute the scalar coordinates $c_1$ and $c_2$ by taking dot products with the unit basis vectors:$c_1 = \mathbf{q}_1^T \mathbf{x} = \begin{bmatrix} \frac{1}{\sqrt{2}} & \frac{1}{\sqrt{2}} \end{bmatrix} \begin{bmatrix} 5 \\ 1 \end{bmatrix} = \frac{6}{\sqrt{2}}$$c_2 = \mathbf{q}_2^T \mathbf{x} = \begin{bmatrix} -\frac{1}{\sqrt{2}} & \frac{1}{\sqrt{2}} \end{bmatrix} \begin{bmatrix} 5 \\ 1 \end{bmatrix} = \frac{-4}{\sqrt{2}}$\\So, in the eigenvector basis, $\mathbf{x}$ is represented as:$$\mathbf{x} = \left(\frac{6}{\sqrt{2}}\right) \mathbf{q}_1 + \left(\frac{-4}{\sqrt{2}}\right) \mathbf{q}_2$$
	Scale Each Component by Eigenvalues: Multiply each coordinate by its respective eigenvalue $\lambda_i$:\\
	Along $\mathbf{q}_1$: $\quad \lambda_1 \cdot c_1 = 4 \cdot \left(\frac{6}{\sqrt{2}}\right) = \frac{24}{\sqrt{2}}$\\
	Along $\mathbf{q}_2$: $\quad \lambda_2 \cdot c_2 = 2 \cdot \left(\frac{-4}{\sqrt{2}}\right) = \frac{-8}{\sqrt{2}}$\\
	Step C: Reconstruct the Final Vector
	$$A\mathbf{x} = \left(\frac{24}{\sqrt{2}}\right) \mathbf{q}_1 + \left(\frac{-8}{\sqrt{2}}\right) \mathbf{q}_2 \implies A\mathbf{x} = \begin{bmatrix} 12 \\ 12 \end{bmatrix} + \begin{bmatrix} 4 \\ -4 \end{bmatrix} = \begin{bmatrix} 16 \\ 8 \end{bmatrix}$$
	
	\takeaway{Instead of viewing matrix multiplication $A\mathbf{x}$ as a single black-box calculation, basis expansion reveals that $A$:
	\begin{itemize}
		\item Breaks $\mathbf{x}$ down into independent components along the principal axes $\mathbf{q}_1$ and $\mathbf{q}_2$.
		\item Scales each component independently by its eigenvalue ($\times 4$ along $\mathbf{q}_1$, $\times 2$ along $\mathbf{q}_2$).
		\item Adds the scaled components together to form the output vector $\begin{bmatrix} 16 \\ 8 \end{bmatrix}$.
		\end{itemize}	}
	
	\subsection{Singular Value Decomposition (SVD):}
	\deff The Singular Value Decomposition factors any real $m \times n$ matrix $A$ into three matrices:$$A = U \Sigma V^T$$
	\begin{itemize}
		\item $U$ ($m \times m$): An orthogonal matrix containing the left singular vectors (eigenvectors of $A A^T$).
		\item $\Sigma$ ($m \times n$): A diagonal matrix containing the singular values $\sigma_i = \sqrt{\lambda_i}$, sorted in descending order.
		\item $V$ ($n \times n$): An orthogonal matrix containing the right singular vectors (eigenvectors of $A^T A$).
	\end{itemize}
	\example Find the SVD of the $2 \times 2$ matrix:$$A = \begin{bmatrix} 3 & 0 \\ 0 & -2 \end{bmatrix}$$
	\solution
	$$A^T A = \begin{bmatrix} 3 & 0 \\ 0 & -2 \end{bmatrix} \begin{bmatrix} 3 & 0 \\ 0 & -2 \end{bmatrix} = \begin{bmatrix} 9 & 0 \\ 0 & 4 \end{bmatrix}$$
	Eigenvalues of $A^T A$: $$\lambda_1 = 9, \quad \lambda_2 = 4$$
	Eigenvectors of $A^T A$:\\
	For $\lambda_1 = 9$: $\mathbf{v}_1 = \begin{bmatrix} 1 \\ 0 \end{bmatrix}$, 
	for $\lambda_2 = 4$: $\mathbf{v}_2 = \begin{bmatrix} 0 \\ 1 \end{bmatrix}$
	$$\Rightarrow V = \begin{bmatrix} 1 & 0 \\ 0 & 1 \end{bmatrix} \implies V^T = \begin{bmatrix} 1 & 0 \\ 0 & 1 \end{bmatrix}$$
	Calculate Singular Values ($\Sigma$)The singular values $\sigma_i$ are the positive square roots of the eigenvalues $\lambda_i$:$$\sigma_1 = \sqrt{9} = 3$$$$\sigma_2 = \sqrt{4} = 2$$Construct the matrix $\Sigma$:$$\Sigma = \begin{bmatrix} \sigma_1 & 0 \\ 0 & \sigma_2 \end{bmatrix} = \begin{bmatrix} 3 & 0 \\ 0 & 2 \end{bmatrix}$$
	Compute Left Singular Vectors ($U$)The columns of $U$ are calculated using the formula:$$\mathbf{u}_i = \frac{1}{\sigma_i} A \mathbf{v}_i$$First Column ($\mathbf{u}_1$):$$\mathbf{u}_1 = \frac{1}{3} \begin{bmatrix} 3 & 0 \\ 0 & -2 \end{bmatrix} \begin{bmatrix} 1 \\ 0 \end{bmatrix} = \frac{1}{3} \begin{bmatrix} 3 \\ 0 \end{bmatrix} = \begin{bmatrix} 1 \\ 0 \end{bmatrix}$$
	Similarly $$\mathbf{u}_2 = \frac{1}{2} \begin{bmatrix} 3 & 0 \\ 0 & -2 \end{bmatrix} \begin{bmatrix} 0 \\ 1 \end{bmatrix} = \frac{1}{2} \begin{bmatrix} 0 \\ -2 \end{bmatrix} = \begin{bmatrix} 0 \\ -1 \end{bmatrix}$$
	$$U = \begin{bmatrix} 1 & 0 \\ 0 & -1 \end{bmatrix}$$
	Thus
	$$A= U \Sigma V^T = \begin{bmatrix} 1 & 0 \\ 0 & -1 \end{bmatrix} \begin{bmatrix} 3 & 0 \\ 0 & 2 \end{bmatrix} \begin{bmatrix} 1 & 0 \\ 0 & 1 \end{bmatrix}$$
	\example
	Let's take the rank-3 matrix:$$A = \begin{bmatrix} 1 & 2 & 3 \\ 1 & 1 & 2 \\ 2 & 1 & 1 \end{bmatrix}$$
	Since all rows are linearly independent ($\det(A) = 2 \neq 0$), the matrix has full rank (Rank = 3). It has 3 non-zero singular values.
	First, calculate $A^T A$:$$A^T A = \begin{bmatrix} 1 & 1 & 2 \\ 2 & 1 & 1 \\ 3 & 2 & 1 \end{bmatrix} \begin{bmatrix} 1 & 2 & 3 \\ 1 & 1 & 2 \\ 2 & 1 & 1 \end{bmatrix} = \begin{bmatrix} 6 & 5 & 7 \\ 5 & 6 & 9 \\ 7 & 9 & 14 \end{bmatrix}$$
	$$\lambda_1 \approx 23.518, \quad \lambda_2 \approx 2.115, \quad \lambda_3 \approx 0.367$$
	$$\mathbf{v}_1 \approx \begin{bmatrix} 0.407 \\ 0.498 \\ 0.766 \end{bmatrix}, \quad \mathbf{v}_2 \approx \begin{bmatrix} -0.840 \\ -0.222 \\ 0.495 \end{bmatrix}, \quad \mathbf{v}_3 \approx \begin{bmatrix} 0.359 \\ -0.838 \\ 0.313 \end{bmatrix}$$
	$$V \approx \begin{bmatrix} 0.407 & -0.840 & 0.359 \\ 0.498 & -0.222 & -0.838 \\ 0.766 & 0.495 & 0.313 \end{bmatrix}$$
	Singular Values ($\Sigma$)
	$$\sigma_1 = \sqrt{23.518} \approx 4.850, \sigma_2 = \sqrt{2.115} \approx 1.454, \sigma_3 = \sqrt{0.367} \approx 0.606$$ 
	Thus $$\Sigma \approx \begin{bmatrix} 4.850 & 0 & 0 \\ 0 & 1.454 & 0 \\ 0 & 0 & 0.606 \end{bmatrix}$$
	 Left Singular Vectors ($U$
	 $$\mathbf{u}_i = \frac{1}{\sigma_i} A \mathbf{v}_i$$
	 $$U \approx \begin{bmatrix} 0.763 & 0.138 & -0.624 \\ 0.502 & -0.051 & 0.243 \\ 0.428 & -0.968 & 0.318 \end{bmatrix}$$
	 Hence
	 $$A = U \Sigma V^T  \approx \begin{bmatrix} 0.763 & 0.138 & -0.624 \\ 0.502 & -0.051 & 0.243 \\ 0.428 & -0.968 & 0.318 \end{bmatrix} \begin{bmatrix} 4.850 & 0 & 0 \\ 0 & 1.454 & 0 \\ 0 & 0 & 0.606 \end{bmatrix} \begin{bmatrix} 0.407 & 0.498 & 0.766 \\ -0.840 & -0.222 & 0.495 \\ 0.359 & -0.838 & 0.313 \end{bmatrix}$$
	 
	 
	 \example  Consider the matrix of rank 1: $$A=\begin{bmatrix} 1 & 2 & 3 & 4 \\ -1 & -2 & -3 & -4 \\ 2 & 4 & 6 & 8 \\ 10 & 20 & 30 & 40 \end{bmatrix}$$
	 The above matrix can be expressed as:
	$$ A= {\color{blue}\begin{bmatrix} 1 \\ -1 \\ 2 \\ 10 \end{bmatrix}} {\color{green!50!black}\begin{bmatrix} 1 & 2 & 3 & 4 \end{bmatrix}}$$
	

	Similarly for the second example of rank 3 matrix, the given matrix $A$ can be expressed as:
	 \begin{align*}
	 	A \approx {\color{red}4.850} {\color{blue}\begin{bmatrix} 0.763 \\ 0.502 \\ 0.428 \end{bmatrix}} {\color{green!50!black}\begin{bmatrix} 0.407 & 0.498 & 0.766 \end{bmatrix}} + {\color{red}1.454} {\color{blue}\begin{bmatrix} 0.138 \\ -0.051 \\ -0.968 \end{bmatrix}} {\color{green!50!black}\begin{bmatrix} -0.840 & -0.222 & 0.495 \end{bmatrix}} \\
	 	+ {\color{red}0.606} {\color{blue}\begin{bmatrix} -0.624 \\ 0.243 \\ 0.318 \end{bmatrix}} {\color{green!50!black}\begin{bmatrix} 0.359 & -0.838 & 0.313 \end{bmatrix}}
	 \end{align*}
	\takeaway{Taking only the first two singular values ($\sigma_1$ and $\sigma_2$) discards the third outer product component, yielding the optimal Rank-2 approximation ($A_2$) of the original matrix $A$:
	$$A = \begin{bmatrix} 1 & 2 & 3 \\ 1 & 1 & 2 \\ 2 & 1 & 1 \end{bmatrix} \approx A_2 = \mathbf{\begin{bmatrix} 1.338 & 1.798 & 2.934 \\ 1.053 & 1.228 & 1.828 \\ 2.027 & 1.347 & 0.893 \end{bmatrix}}$$}
	 
	 \takeaway{An $m \times n$ matrix requires $m \cdot n$ stored values. Representing it as a rank-1 outer product of two vectors requires only $m + n$ values. For a $1000 \times 1000$ matrix, storage drops from $1,000,000$ numbers down to just $2,000$—a 99.8$\%$ reduction.
	 \par
	 Tiny Matrices ($3 \times 3$): Full SVD adds mathematical overhead because storing the three matrices ($U, \Sigma, V^T$) takes 21 values. You only save space on a $3 \times 3$ matrix if you truncate to Rank-1 (storing 7 values).
	 \par
	\textbf{ Where SVD Truly Shines (Large Matrices):} For an $m \times n$ matrix, Rank-$k$ SVD requires $k(1 + m + n)$ values instead of $m \cdot n$. For a $1000 \times 1000$ matrix ($1,000,000$ original values), keeping $k = 20$ singular components requires only $20(1 + 1000 + 1000) = \mathbf{40,020\text{ values}}$, resulting in a 96$\%$ storage reduction.}
	 
	\example
	Imagine a high resolution $1024 \times 1024$ image (See Figure \ref{fig:image_comparison}). Mathematically, it is a matrix $A$ containing 1,048,576 individual pixel values. Storing, transmitting, or processing every single pixel individually is computationally expensive. However, real world images aren't random noise neighboring pixels share patterns, shapes, and gradients. Singular Value Decomposition (SVD) exploits this redundancy by factorizing the original matrix into its most fundamental building blocks:$$A = U \Sigma V^T$$Instead of storing all 1+ million numbers, SVD ranks these patterns by importance (their singular values) and allows us to construct a low-rank approximation $A_r$:$$A_r = U_r \Sigma_r V_r^T$$
	\begin{figure}[H]
		\centering
		
		% First row
		\begin{minipage}{0.3\textwidth}
			\centering
			\includegraphics[width=\linewidth]{Original}
		\end{minipage}
		\hspace{0.01\textwidth}
		\begin{minipage}{0.3\textwidth}
			\centering
			\includegraphics[width=\linewidth]{25 compressed}
		\end{minipage}
		
		\vspace{0.5cm}
		
		% Second row
		\begin{minipage}{0.3\textwidth}
			\centering
			\includegraphics[width=\linewidth]{50 compressed}
		\end{minipage}
		\hspace{0.01\textwidth}
		\begin{minipage}{0.3\textwidth}
			\centering
			\includegraphics[width=\linewidth]{75 compressed}
		\end{minipage}
		
		\caption{Comparison of the original $1024\times1024$ image 
			(1.98 MB) and the compressed $1024\times1024$ images 
			obtained using Compression Rank 25 (113 KB), 50 (135 KB) and 75 (148 KB).}
		\label{fig:image_comparison}
	\end{figure}
	
	By retaining only the top $r$ most dominant singular values and their corresponding singular vectors:
	\begin{itemize}
		\item Original Matrix Storage: $1024 \times 1024 = \mathbf{1,048,576\text{ values}}$
		\item SVD Rank-$r$ Storage: $1024r + r + 1024r = \mathbf{2049r\text{ values}}$
	\end{itemize}
		If we set $r = 75$:$$\text{Storage Required} = 2049 \times 75 = \mathbf{153,675\text{ values}}$$The Result: We preserve over 85$\%$ to 90$\%$ of the visual quality while storing less than 15$\%$ of the original data.
		
	\takeaway{A massive matrix may look overwhelmingly complex, but SVD reveals that most of its critical information is concentrated along a remarkably small set of dominant directions.
	}
	
	\subsection{Principal Component Analysis (PCA)}
	PCA transforms a dataset containing $p$ potentially correlated variables into a smaller set of $k$ ($k \le p$) uncorrelated, orthogonal variables called Principal Components (PCs).\\
	These principal components are constructed as linear combinations of the original variables in such a way that:
	\begin{itemize}
		\item 	The first principal component captures the maximum possible variance present in the dataset.
		
		\item Each subsequent component captures the maximum remaining variance subject to the constraint that it is orthogonal (uncorrelated) to all preceding components.
	\end{itemize}

	
	Mathematically, PCA rotates the original coordinate axes to align with the directions of maximum variance in the data feature space.	
\end{itemize}

\textbf{Procedure: } For a given matrix $A_{p\times n}$, i.e., $p$ features and $n$ observations\\
\begin{itemize}
	\item Step 1: Compute the Mean Vector ($\mathbf{\bar{x}}$), $\bar{x}_i$ = Mean of feature $i$ calculated across all $n$ samples:$$\bar{x}_i = \frac{1}{n} \sum_{k=1}^{n} x_{i, k}$$$X_{i, k} = (x_{i, k} - \bar{x}_i)$ = Mean-centered scalar value for feature $i$ in sample $k$.\\
	Construct the $p \times n$ mean-centered data matrix $\bar{A}$:$$\bar{A} = \begin{bmatrix} X_{1,1} & X_{1,2} & X_{1,3} & \dots & X_{1,n} \\ X_{2,1} & X_{2,2} & X_{2,3} & \dots & X_{2,n} \\ \vdots & \vdots & \vdots & \ddots & \vdots \\ X_{p,1} & X_{p,2} & X_{p,3} & \dots & X_{p,n} \end{bmatrix}_{p \times n}$$
	\item Step 2: Compute the Covariance Matrix ($\Sigma$):\\
		\textbf{Method 01:} $$\boldsymbol{\Sigma}_{\text{cov}} = \begin{bmatrix} \boldsymbol{\Sigma}_{1,1} & \boldsymbol{\Sigma}_{1,2} & \boldsymbol{\Sigma}_{1,3} & \dots & \boldsymbol{\Sigma}_{1,p} \\ \boldsymbol{\Sigma}_{2,1} & \boldsymbol{\Sigma}_{2,2} & \boldsymbol{\Sigma}_{2,3} & \dots & \boldsymbol{\Sigma}_{2,p} \\ \boldsymbol{\Sigma}_{3,1} & \boldsymbol{\Sigma}_{3,2} & \boldsymbol{\Sigma}_{3,3} & \dots & \boldsymbol{\Sigma}_{3,p} \\ \vdots & \vdots & \vdots & \ddots & \vdots \\ \boldsymbol{\Sigma}_{p,1} & \boldsymbol{\Sigma}_{p,2} & \boldsymbol{\Sigma}_{p,3} & \dots & \boldsymbol{\Sigma}_{p,p} \end{bmatrix}$$
	 Diagonal Elements ($i = j$): Represents the variance of feature $i$:$$\boldsymbol{\Sigma}_{i, i} = \text{Var}(x_i) = \frac{1}{n - 1} \sum_{k=1}^{n} X_{i, k}^2$$Off-Diagonal Elements ($i \neq j$): Represents the covariance between feature $i$ and feature $j$:$$\boldsymbol{\Sigma}_{i, j} = \text{Cov}(x_i, x_j) = \frac{1}{n - 1} \sum_{k=1}^{n} X_{i, k} X_{j, k}$$
	\textbf{Method 02: } 
	Find 
	$$AA^T = \begin{bmatrix} X_{1,1} & \dots & X_{1,n} \\ X_{2,1} & \dots & X_{2,n} \\ \vdots & \ddots & \vdots \\ X_{p,1} & \dots & X_{p,n} \end{bmatrix} \begin{bmatrix} X_{1,1} & X_{2,1} & \dots & X_{p,1} \\ X_{1,2} & X_{2,2} & \dots & X_{p,2} \\ \vdots & \vdots & \ddots & \vdots \\ X_{1,n} & X_{2,n} & \dots & X_{p,n} \end{bmatrix}_{n \times p}$$				
	Scaling $AA^T$ by $\frac{1}{n - 1}$ gives the $p \times p$ sample covariance matrix:$$\boldsymbol{\Sigma}_{\text{cov}} = \frac{1}{n - 1} AA^T$$
	\item Step 3: Calculate the Eigenvalues ($\lambda$)
	\item Step 4: Calculate Normalized Eigenvectors ($\mathbf{v}_1, \mathbf{v}_2$)
	\item Step 5: Percentage of Variance Explained
	\item Step 6: Compute Principal Component Scores (Transformed Data)
\end{itemize} 

\example Consider 4 documents:
\begin{itemize}
	\item "king crown realm" $\rightarrow$ Doc 0: King context
	\item "queen throne realm" $\rightarrow$ Doc 1: Queen context
	\item "king queen throne crown" $\rightarrow$ Doc 2: Shared context link
	\item "dog puppy canine" $\rightarrow$ Doc 3: Pet context A
	\item "cat kitten feline" $\rightarrow$ Doc 4: Pet context B
	\item "dog cat puppy kitten" $\rightarrow$ Doc 5: Shared pet link
\end{itemize}
If we want to print the vector for Doc 0 ("king crown realm"), it would contain 11 numbers. Each number would represent the count of each of the 11 vocabulary words inside that specific document called as One-Hot Vectors. The word 'king' is purely as an isolated word (a strict one-hot vector) rather than by its document context, it would indeed be an 11-dimensional vector with a single 1 at the 6th position: [0, 0, 0, 0, 0, 1, 0, 0, 0, 0, 0].\\
There is a famous rule in linguistics formulated by John Rupert Firth:

"You shall know a word by the company it keeps."

Instead of treating words as isolated standard basis vectors, we build a Term-Document / Co-Occurrence Matrix $A$. Each row corresponds to a word, and each column corresponds to a document or context.
For example 
\begin{lstlisting}[language=Python, caption=Rows from Co-Occurrence Matrix]
	'king' context vector  : [1 0 1 0 0 0]
	'queen' context vector : [0 1 1 0 0 0]
	'dog' context vector : [0 0 0 1 0 1]
	#Document wise:
	["canine", "cat", "crown", "dog","feline", "king", "kitten", "puppy", "queen", "realm", "throne"]
	Doc 0 ->[0, 0, 1, 0, 0, 1, 0, 0, 0, 1, 0]
	Doc 1 ->[0, 0, 0, 0, 0, 0, 0, 0, 1, 1, 1]
	Doc 2 ->[0, 0, 1, 0, 0, 1, 0, 0, 1, 0, 1]
	Doc 3 ->[1, 0, 0, 1, 0, 0, 0, 1, 0, 0, 0]
	Doc 4 ->[0, 1, 0, 0, 1, 0, 1, 0, 0, 0, 0]
	Doc 5 ->[0, 1, 0, 1, 0, 0, 1, 1, 0, 0, 0]
\end{lstlisting}
SVD looks at these 6-dimensional document occurrence vectors and notices that Doc 0, Doc 1, and Doc 2 share words like "realm", "throne", and "crown". It compresses the 6 document dimensions into a small number of latent topic axes (like Royalty vs. Pets), where 'king' and 'queen' get almost identical coordinates.
\begin{lstlisting}[language=Python, caption=Word Coordinates in 2D Latent Feature Space:]
	        Topic Axis 1  Topic Axis 2
	canine        -0.000         0.180
	cat            0.000         0.484
	crown          0.462        -0.000
	dog            0.000         0.484
	feline         0.000         0.180
	king           0.462        -0.000
	kitten         0.000         0.484
	puppy          0.000         0.484
	queen          0.462        -0.000
	realm          0.383        -0.000
	throne         0.462        -0.000
\end{lstlisting}

	\begin{lstlisting}[language=Python, caption=Python Code]
	import numpy as np
	import pandas as pd
	from sklearn.feature_extraction.text import CountVectorizer
	from sklearn.decomposition import TruncatedSVD
	from sklearn.metrics.pairwise import cosine_similarity
	
	# 1. SAMPLE TEXT CORPUS
	# Notice how 'king' and 'queen' never appear in the same document directly,
	# but both co-occur with shared words ('crown', 'throne', 'realm').
	documents = [
	"king crown realm",          # Doc 0: King context
	"queen throne realm",        # Doc 1: Queen context
	"king queen throne crown",   # Doc 2: Shared context link
	"dog puppy canine",          # Doc 3: Pet context A
	"cat kitten feline",         # Doc 4: Pet context B
	"dog cat puppy kitten"       # Doc 5: Shared pet link
	]
	
	# 2. RAW HIGH-DIMENSIONAL ONE-HOT SPACE
	vectorizer = CountVectorizer()
	X_raw = vectorizer.fit_transform(documents).toarray()
	words = vectorizer.get_feature_names_out()
	
	# Extract individual one-hot-like basis vectors for words
	king_idx = list(words).index('king')
	queen_idx = list(words).index('queen')
	dog_idx = list(words).index('dog')
	
	king_raw = X_raw[:, king_idx]
	queen_raw = X_raw[:, queen_idx]
	dog_raw = X_raw[:,dog_idx]
	
	# Calculate cosine similarity BEFORE SVD
	sim_before_king_queen = cosine_similarity([king_raw], [queen_raw])[0][0]
	sim_before_dog_queen = cosine_similarity([dog_raw], [queen_raw])[0][0]
	
	
	print("================ 1. BEFORE SVD (RAW SPACE) ================")
	print(f"Vocabulary ({len(words)} dimensions): {list(words)}")
	print(f"'king' context vector  : {king_raw}")
	print(f"'queen' context vector : {queen_raw}")
	print(f"'dog' context vector : {dog_raw}")
	print(f"Cosine Similarity between King and Queen BEFORE SVD: {sim_before_king_queen:.2f}")
	print(f"Cosine Similarity between Dog and Queen BEFORE SVD: {sim_before_dog_queen:.2f}")
	
	# 3. REDUCE DIMENSIONS VIA SVD (Change Basis to Latent Topics)
	svd = TruncatedSVD(n_components=2, random_state=42)
	X_svd = svd.fit_transform(X_raw)
	#Under the hood, it operates on the term co-occurrence matrix ($X_{\text{raw}}^T X_{\text{raw}}$) to calculate singular values and vectors
		
	# Word coordinates in the new 2D latent basis space (U * Sigma)
	word_vectors_svd = svd.components_.T
	#In scikit-learn, the .components_ attribute explicitly stores the right singular vectors, $V^T$.
	
	king_svd = word_vectors_svd[king_idx]
	queen_svd = word_vectors_svd[queen_idx]
	dog_svd = word_vectors_svd[dog_idx]
	
	# Calculate cosine similarity AFTER SVD
	sim_after_king_queen = cosine_similarity([king_svd], [queen_svd])[0][0]
	sim_after_dog_queen = cosine_similarity([dog_svd], [queen_svd])[0][0]
	
	
	
	print("\n================ 2. AFTER SVD (NEW LATENT FEATURE BASIS) ================")
	df_word_coords = pd.DataFrame(
	data=word_vectors_svd,
	index=words,
	columns=["Topic Axis 1", "Topic Axis 2"]
	).round(3)
	
	print("Word Coordinates in 2D Latent Feature Space:")
	print(df_word_coords)
	
	print(f"\nCompressed 'king' vector  : {np.round(king_svd, 3)}")
	print(f"Compressed 'queen' vector : {np.round(queen_svd, 3)}")
	print(f"Compressed 'dog' vector : {np.round(dog_svd, 3)}")
	print(f"Cosine Similarity between King and Queen AFTER SVD : {sim_after_king_queen:.2f}")
	print(f"Cosine Similarity between Dog and Queen  AFTER SVD : {sim_after_dog_queen:.2f}")
\end{lstlisting}

\begin{lstlisting}[language=Python, caption=Output]
================ 1. BEFORE SVD (RAW SPACE) ================
Vocabulary (11 dimensions): ['canine', 'cat', 'crown', 'dog', 'feline', 'king', 'kitten', 'puppy', 'queen', 'realm', 'throne']
'king' context vector  : [1 0 1 0 0 0]
'queen' context vector : [0 1 1 0 0 0]
'dog' context vector : [0 0 0 1 0 1]
Cosine Similarity between King and Queen BEFORE SVD: 0.50
Cosine Similarity between Dog and Queen BEFORE SVD: 0.00

================ 2. AFTER SVD (NEW LATENT FEATURE BASIS) ================
Word Coordinates in 2D Latent Feature Space:
 Topic Axis 1  Topic Axis 2
canine        -0.000         0.180
cat            0.000         0.484
crown          0.462        -0.000
dog            0.000         0.484
feline         0.000         0.180
king           0.462        -0.000
kitten         0.000         0.484
puppy          0.000         0.484
queen          0.462        -0.000
realm          0.383        -0.000
throne         0.462        -0.000

Compressed 'king' vector  : [ 0.462 -0.   ]
Compressed 'queen' vector : [ 0.462 -0.   ]
Compressed 'dog' vector : [0.    0.484]
Cosine Similarity between King and Queen AFTER SVD : 1.00
Cosine Similarity between Dog and Queen  AFTER SVD : -0.00}
\end{lstlisting}

\begin{lstlisting}[language=Python, caption=Output]
	import gensim.downloader as api
	from gensim.models import Word2Vec
	from sklearn.decomposition import PCA
	import matplotlib.pyplot as plt
	
	# 1. Download and train on Wikipedia sample data (contains math, science, and everyday words)
	print("Loading text8 dataset...")
	dataset = api.load('text8')
	
	print("Training Word2Vec...")
	model = Word2Vec(sentences=dataset, vector_size=100, window=5, min_count=5, workers=1)
	
	# 2. Define your target words
	target_words = [
	'line', 'triangle', 'curve', 'calculus', 'algebra', 'mathematics', 'math', # Math
	'love', 'happy', 'sad', 'hope',                                    # Emotions
	'keyboard', 'pen', 'car', 'computer'                               # Objects/Actions
	]
	
	# 3. Filter and Plot
	plot_words = [w for w in target_words if w in model.wv]
	plot_vectors = [model.wv[w] for w in plot_words]
	
	pca = PCA(n_components=2)
	vectors_2d = pca.fit_transform(plot_vectors)
	
	plt.figure(figsize=(10, 7))
	plt.scatter(vectors_2d[:, 0], vectors_2d[:, 1], color='dodgerblue', edgecolors='k', s=80)
	
	for i, word in enumerate(plot_words):
	plt.annotate(word, xy=(vectors_2d[i, 0], vectors_2d[i, 1]), xytext=(6, 3), textcoords='offset points')
	
	plt.title("2D PCA Projection of Word Vectors (Wikipedia Dataset)")
	plt.grid(True, linestyle='--', alpha=0.5)
	plt.show()
	
	# 4. Find cosine similarity between two words in the space
	word1 = 'mathematics'
	word2 = 'math'
	print("\nSimilarity ('{word1}', '{word2}'):", model.wv.similarity(word1, word2))
\end{lstlisting}

\begin{figure}[H]
	\centering
	\includegraphics[width=0.5\linewidth]{PCA_Output01}
	\caption{}
	\label{fig:pcaoutput01}
\end{figure}














\newpage
\textbf{References:}

\begin{itemize}
	\item C. J. Reynolds, ``I Built an LLM From Scratch,'' YouTube, Syntax, [Online]. 
	Available: \url{https://youtu.be/YmLp8qe87A0?si=1od0Qcs9OAV7nxhi}.
	
	\item S. Brunton, ``SVD: Image Compression [Python],'' YouTube, [Online]. 
	Available: \url{https://youtu.be/H7qMMudo3e8}.
\end{itemize}




\begin{lstlisting}[language=Python, caption=Python Code]
	import gensim.downloader as api
	from gensim.models import Word2Vec
	from sklearn.decomposition import PCA
	import matplotlib.pyplot as plt
	
	# 1. Download and train on Wikipedia sample data (contains math, science, and everyday words)
	print("Loading text8 dataset...")
	dataset = api.load('text8')
	
	print("Training Word2Vec...")
	model = Word2Vec(sentences=dataset, vector_size=100, window=5, min_count=5, workers=1)
	
	# 2. Define your target words
	target_words = [
	'line', 'triangle', 'curve', 'calculus', 'algebra', 'mathematics', # Math
	'love', 'happy', 'sad', 'hope',                                    # Emotions
	'keyboard', 'pen', 'car', 'computer'                               # Objects/Actions
	]
	
	# 3. Filter and Plot
	plot_words = [w for w in target_words if w in model.wv]
	plot_vectors = [model.wv[w] for w in plot_words]
	
	pca = PCA(n_components=2)
	vectors_2d = pca.fit_transform(plot_vectors)
	
	plt.figure(figsize=(10, 7))
	plt.scatter(vectors_2d[:, 0], vectors_2d[:, 1], color='dodgerblue', edgecolors='k', s=80)
	
	for i, word in enumerate(plot_words):
	plt.annotate(word, xy=(vectors_2d[i, 0], vectors_2d[i, 1]), xytext=(6, 3), textcoords='offset points')
	
	plt.title("2D PCA Projection of Word Vectors (Wikipedia Dataset)")
	plt.grid(True, linestyle='--', alpha=0.5)
	plt.show()
\end{lstlisting}
\end{document}





Word2Vec doesn't manually select or label features; instead, it discovers them automatically through an optimization task using backpropagation and gradient descent.Step 1: The Prediction TaskWord2Vec creates a "fake" classification game to force the network to learn. In the Skip-gram architecture, the task is: Given a target word $w_t$, predict which context words $w_c$ are likely to appear nearby.Step 2: Vector Math (Dot Product)To check how likely word $B$ is to appear near word $A$, the model calculates the dot product of their feature vectors:$$\text{Similarity Score} = \vec{v}_A \cdot \vec{v}_B$$A higher dot product means the model thinks the two words belong in the same context.Initially, because vector features start as random numbers near zero, every dot product score is practically zero.Step 3: Calculating Error (Loss Function)The model evaluates its guess using Negative Sampling:Positive Pair: Words actually seen together in text (e.g., 'dog' and 'barks'). The target score is $1.0$.Negative Pairs: Words randomly sampled that don't appear together (e.g., 'dog' and 'algebra'). The target score is $0.0$.Step 4: Adjusting Weights via Gradient DescentIf the model sees 'dog' and 'barks' together, but their current vectors produce a low dot product, it calculates an error gradient. Through backpropagation, it slightly nudges each of the 100 vector features:Pulling together: Feature values for 'dog' and 'barks' are adjusted to move closer in 100-dimensional space.Pushing apart: Feature values for 'dog' and negative words (like 'algebra') are adjusted to move further apart.How Meaning EmergesAfter repeating this process millions of times across a dataset:Words that consistently share context words (e.g., 'dog' and 'puppy' both appearing near 'fur', 'vet', and 'cute') get pulled toward the same coordinate regions.The 100 dimensions naturally align into abstract axes representing grammar, concepts, topics, and relationships—without any human needing to name or define them.

